\documentclass[11pt,oneside,openany]{book}

\usepackage[a4paper,left=18mm,right=18mm,top=22mm,bottom=19mm,headheight=21pt]{geometry}
\usepackage{fontspec}
\setmainfont{Arial}
\setsansfont{Arial}
\setmonofont{Consolas}[Scale=0.88]

\usepackage[table]{xcolor}
\definecolor{PageBg}{HTML}{0B1F33}
\definecolor{PageBgSoft}{HTML}{102A43}
\definecolor{Panel}{HTML}{153B4A}
\definecolor{PanelAlt}{HTML}{16324A}
\definecolor{BodyText}{HTML}{F4F8FB}
\definecolor{MutedText}{HTML}{B8CDD9}
\definecolor{Accent}{HTML}{34D1BF}
\definecolor{AccentTwo}{HTML}{F6AE2D}
\definecolor{RuleColor}{HTML}{315872}
\pagecolor{PageBg}
\color{BodyText}

\usepackage{amsmath,amssymb}
\usepackage{graphicx}
\usepackage{longtable,booktabs,array,calc}
\newcounter{none}
\usepackage{etoolbox}
\usepackage{enumitem}
\usepackage{fancyvrb}
\usepackage{fvextra}
\usepackage{framed}
\usepackage{titlesec}
\usepackage{fancyhdr}
\usepackage{tikz}
\usetikzlibrary{arrows.meta,positioning}
\usepackage[unicode,colorlinks=true,linkcolor=Accent,urlcolor=Accent,citecolor=Accent]{hyperref}
\usepackage{bookmark}
\hypersetup{
  pdftitle={Cẩm nang Thái Kiệt - SafeRoom IoT Box},
  pdfauthor={SafeRoom IoT Box},
  pdfsubject={CB-03, NC-01, NC-03, NC-05}
}
\renewcommand{\contentsname}{Mục lục}

\setlength{\parindent}{0pt}
\setlength{\parskip}{4pt plus 1pt minus 1pt}
\setlength{\emergencystretch}{3em}
\renewcommand{\arraystretch}{1.25}
\let\originaltexttt\texttt
\renewcommand{\texttt}[1]{{\ttfamily\small #1}}
\setcounter{secnumdepth}{-1}
\setcounter{tocdepth}{1}

\titleformat{\chapter}[display]
  {\sffamily\bfseries\color{BodyText}}
  {\color{Accent}\large PHẦN}
  {4pt}
  {\Huge}
  [\vspace{4pt}{\color{AccentTwo}\titlerule[1.2pt]}]
\titlespacing*{\chapter}{0pt}{-8pt}{22pt}
\titleformat{\section}{\sffamily\bfseries\Large\color{Accent}}{}{0pt}{}
\titleformat{\subsection}{\sffamily\bfseries\large\color{AccentTwo}}{}{0pt}{}
\titleformat{\subsubsection}{\sffamily\bfseries\normalsize\color{Accent}}{}{0pt}{}
\titlespacing*{\section}{0pt}{15pt}{5pt}
\titlespacing*{\subsection}{0pt}{11pt}{4pt}
\titlespacing*{\subsubsection}{0pt}{9pt}{3pt}

\pagestyle{fancy}
\fancyhf{}
\fancyhead[L]{\sffamily\bfseries\scriptsize\color{MutedText} CẨM NANG THÁI KIỆT}
\fancyhead[R]{\sffamily\scriptsize\color{MutedText} SAFEROOM IOT BOX}
\fancyfoot[L]{\sffamily\scriptsize\color{MutedText} CB-03 \textbullet{} NC-01 \textbullet{} NC-03 \textbullet{} NC-05}
\fancyfoot[R]{\sffamily\bfseries\color{Accent}\thepage}
\renewcommand{\headrulewidth}{0.35pt}
\renewcommand{\headrule}{\hbox to\headwidth{\color{RuleColor}\leaders\hrule height \headrulewidth\hfill}}
\renewcommand{\footrulewidth}{0pt}
\fancypagestyle{plain}{\pagestyle{fancy}}

\setlist[itemize]{leftmargin=5mm,itemsep=2pt,topsep=3pt,label=\textcolor{Accent}{\textbullet}}
\setlist[enumerate]{leftmargin=6mm,itemsep=2pt,topsep=3pt,label=\textcolor{Accent}{\arabic*.}}

\AtBeginEnvironment{longtable}{\small\color{BodyText}\rowcolors{2}{PageBgSoft}{PanelAlt}}
\AtBeginEnvironment{tabular}{\small\color{BodyText}}
\arrayrulecolor{RuleColor}
\colorlet{tablehead}{Accent}
\let\oldtoprule\toprule
\renewcommand{\toprule}{\arrayrulecolor{Accent}\oldtoprule\arrayrulecolor{RuleColor}}

\usepackage{color}
\usepackage{fancyvrb}
\newcommand{\VerbBar}{|}
\newcommand{\VERB}{\Verb[commandchars=\\\{\}]}
\DefineVerbatimEnvironment{Highlighting}{Verbatim}{commandchars=\\\{\}}
% Add ',fontsize=\small' for more characters per line
\usepackage{framed}
\definecolor{shadecolor}{RGB}{35,38,41}
\newenvironment{Shaded}{\begin{snugshade}}{\end{snugshade}}
\newcommand{\AlertTok}[1]{\textcolor[rgb]{0.58,0.85,0.30}{\textbf{\colorbox[rgb]{0.30,0.12,0.14}{#1}}}}
\newcommand{\AnnotationTok}[1]{\textcolor[rgb]{0.25,0.50,0.35}{#1}}
\newcommand{\AttributeTok}[1]{\textcolor[rgb]{0.16,0.50,0.73}{#1}}
\newcommand{\BaseNTok}[1]{\textcolor[rgb]{0.96,0.45,0.00}{#1}}
\newcommand{\BuiltInTok}[1]{\textcolor[rgb]{0.50,0.55,0.55}{#1}}
\newcommand{\CharTok}[1]{\textcolor[rgb]{0.24,0.68,0.91}{#1}}
\newcommand{\CommentTok}[1]{\textcolor[rgb]{0.48,0.49,0.49}{#1}}
\newcommand{\CommentVarTok}[1]{\textcolor[rgb]{0.50,0.55,0.55}{#1}}
\newcommand{\ConstantTok}[1]{\textcolor[rgb]{0.15,0.68,0.68}{\textbf{#1}}}
\newcommand{\ControlFlowTok}[1]{\textcolor[rgb]{0.99,0.74,0.29}{\textbf{#1}}}
\newcommand{\DataTypeTok}[1]{\textcolor[rgb]{0.16,0.50,0.73}{#1}}
\newcommand{\DecValTok}[1]{\textcolor[rgb]{0.96,0.45,0.00}{#1}}
\newcommand{\DocumentationTok}[1]{\textcolor[rgb]{0.64,0.20,0.25}{#1}}
\newcommand{\ErrorTok}[1]{\textcolor[rgb]{0.85,0.27,0.33}{\underline{#1}}}
\newcommand{\ExtensionTok}[1]{\textcolor[rgb]{0.00,0.60,1.00}{\textbf{#1}}}
\newcommand{\FloatTok}[1]{\textcolor[rgb]{0.96,0.45,0.00}{#1}}
\newcommand{\FunctionTok}[1]{\textcolor[rgb]{0.56,0.27,0.68}{#1}}
\newcommand{\ImportTok}[1]{\textcolor[rgb]{0.15,0.68,0.38}{#1}}
\newcommand{\InformationTok}[1]{\textcolor[rgb]{0.77,0.36,0.00}{#1}}
\newcommand{\KeywordTok}[1]{\textcolor[rgb]{0.81,0.81,0.76}{\textbf{#1}}}
\newcommand{\NormalTok}[1]{\textcolor[rgb]{0.81,0.81,0.76}{#1}}
\newcommand{\OperatorTok}[1]{\textcolor[rgb]{0.81,0.81,0.76}{#1}}
\newcommand{\OtherTok}[1]{\textcolor[rgb]{0.15,0.68,0.38}{#1}}
\newcommand{\PreprocessorTok}[1]{\textcolor[rgb]{0.15,0.68,0.38}{#1}}
\newcommand{\RegionMarkerTok}[1]{\textcolor[rgb]{0.16,0.50,0.73}{\colorbox[rgb]{0.08,0.19,0.26}{#1}}}
\newcommand{\SpecialCharTok}[1]{\textcolor[rgb]{0.24,0.68,0.91}{#1}}
\newcommand{\SpecialStringTok}[1]{\textcolor[rgb]{0.85,0.27,0.33}{#1}}
\newcommand{\StringTok}[1]{\textcolor[rgb]{0.96,0.31,0.31}{#1}}
\newcommand{\VariableTok}[1]{\textcolor[rgb]{0.15,0.68,0.68}{#1}}
\newcommand{\VerbatimStringTok}[1]{\textcolor[rgb]{0.85,0.27,0.33}{#1}}
\newcommand{\WarningTok}[1]{\textcolor[rgb]{0.85,0.27,0.33}{#1}}
\definecolor{shadecolor}{HTML}{183846}
\renewenvironment{Shaded}{%
  \def\FrameCommand{\fboxsep=8pt\colorbox{shadecolor}}%
  \MakeFramed{\advance\hsize-\width\FrameRestore}%
}{\endMakeFramed}
\fvset{breaklines=true,breakanywhere=true,fontsize=\small}

\renewenvironment{quote}{%
  \def\FrameCommand{\hspace{1pt}{\color{Accent}\vrule width 3pt}\hspace{9pt}}%
  \MakeFramed{\advance\hsize-\width\FrameRestore}%
  \color{MutedText}\itshape
}{\endMakeFramed}

\providecommand{\tightlist}{\setlength{\itemsep}{1pt}\setlength{\parskip}{0pt}}
\providecommand{\pandocbounded}[1]{#1}
\providecommand{\passthrough}[1]{#1}
\newcommand{\sectionbreak}{\clearpage}

\begin{document}

\begin{titlepage}
\pagecolor{PageBgSoft}
\thispagestyle{empty}
\begin{tikzpicture}[remember picture,overlay]
  \fill[Accent!45!PageBgSoft] (current page.north east) circle (42mm);
  \fill[Accent!15!PageBgSoft] (current page.south west) circle (34mm);
\end{tikzpicture}
\vspace*{37mm}
\begin{center}
  {\sffamily\bfseries\color{AccentTwo}\large SAFEROOM IOT BOX\par}
  \vspace{8mm}
  {\sffamily\bfseries\color{BodyText}\fontsize{30}{35}\selectfont CẨM NANG THÁI KIỆT\par}
  \vspace{5mm}
  {\sffamily\color{MutedText}\Large Vấn đáp - Q\&A - Code walkthrough - Kịch bản demo\par}
  \vspace{18mm}
  {\color{Accent}\rule{0.68\textwidth}{1.1pt}\par}
  \vspace{8mm}
  {\sffamily\bfseries\color{BodyText}\large CB-03 \quad NC-01 \quad NC-03 \quad NC-05\par}
  \vspace{26mm}
  {\sffamily\color{MutedText}\large
   Điều khiển quạt \textbullet{} Anomaly engine \textbullet{} MongoDB Time Series\par
   Email cảnh báo và throttle 15 phút\par}
\end{center}
\vfill
\begin{center}
  {\sffamily\small\color{MutedText} Tài liệu hợp nhất từ bốn hướng dẫn trong thư mục \texttt{docs}\par}
\end{center}
\end{titlepage}

\pagecolor{PageBg}
\thispagestyle{empty}
\chapter*{Cách dùng cẩm nang}
\addcontentsline{toc}{chapter}{Cách dùng cẩm nang}
Tài liệu được sắp theo bốn nhịp học: hiểu kiến trúc và câu trả lời cốt lõi; luyện 80 câu hỏi; theo dấu code; cuối cùng là chạy demo và xử lý lỗi. Hãy hiểu luồng dữ liệu trước, nhớ contract và giới hạn sau, rồi mới học vị trí file.

\begin{center}
\renewcommand{\arraystretch}{1.45}
\begin{tabular}{>{\bfseries\color{Accent}}p{0.34\textwidth}p{0.56\textwidth}}
Ôn nhanh trước vấn đáp & Đọc Phần I và cheat sheet cuối phần.\\
Luyện phản biện & Đọc Phần II - 80 câu hỏi và trả lời.\\
Chuẩn bị mở code & Đọc Phần III - Code walkthrough.\\
Chuẩn bị demo & Đọc Phần IV - Kịch bản demo.\\
\end{tabular}
\end{center}

\clearpage
\thispagestyle{empty}
\tableofcontents
\clearpage
\setcounter{page}{1}

\chapter{Cẩm nang vấn đáp phần Thái
Kiệt}\label{cux1ea9m-nang-vux1ea5n-ux111uxe1p-phux1ea7n-thuxe1i-kiux1ec7t}

Tài liệu này bám theo proposal của nhóm và phiên bản code hiện tại của
SafeRoom IoT Box. Mục tiêu là giúp Thái Kiệt giải thích đúng phần mình
phụ trách, chỉ được code liên quan, nói rõ giới hạn và đưa ra bằng chứng
kiểm thử.

\section{1. Phạm vi phụ trách}\label{1-phux1ea1m-vi-phux1ee5-truxe1ch}

{\def\LTcaptype{none} % do not increment counter
\begin{longtable}[]{@{}
  >{\raggedright\arraybackslash}p{(\linewidth - 6\tabcolsep) * \real{0.1200}}
  >{\raggedright\arraybackslash}p{(\linewidth - 6\tabcolsep) * \real{0.2600}}
  >{\raggedleft\arraybackslash}p{(\linewidth - 6\tabcolsep) * \real{0.1400}}
  >{\raggedright\arraybackslash}p{(\linewidth - 6\tabcolsep) * \real{0.4300}}@{}}
\toprule\noalign{}
\begin{minipage}[b]{\linewidth}\raggedright
ID
\end{minipage} & \begin{minipage}[b]{\linewidth}\raggedright
Chức năng
\end{minipage} & \begin{minipage}[b]{\linewidth}\raggedleft
Điểm trong proposal
\end{minipage} & \begin{minipage}[b]{\linewidth}\raggedright
Phần cần nắm chắc
\end{minipage} \\
\midrule\noalign{}
\endhead
\bottomrule\noalign{}
\endlastfoot
CB-03 & Điều khiển quạt từ Website & 4.0 & React -\textgreater{} REST
API -\textgreater{} MQTT QoS 1 -\textgreater{} ESP32 -\textgreater{}
relay GPIO 16 \\
NC-01 & Phát hiện dữ liệu bất thường bằng Z-score & 2.0 & Rolling window
N=30, Z-score, hard rules, dedupe, hành động tự động \\
NC-03 & Lưu dữ liệu cảm biến theo thời gian & 1.5 & MongoDB Time Series,
schema, ingestion, retention 30 ngày, historical query \\
NC-05 & Cảnh báo tự động qua Email & 1.0 & Nodemailer, HIGH/CRITICAL,
HTML template, throttle 15 phút \\
\end{longtable}
}

Phần của Kiệt đi xuyên nhiều tầng. Không nên trả lời ``em chỉ làm
backend''. Câu trả lời chính xác hơn:

\begin{quote}
Em phụ trách thiết kế và tích hợp bốn luồng CB-03, NC-01, NC-03 và
NC-05. Trọng tâm code của em nằm ở backend, nhưng CB-03 còn có phần gọi
API ở React và phần nhận lệnh MQTT trên ESP32. Em cũng viết test để
chứng minh các tầng nối với nhau đúng contract.
\end{quote}

\section{2. Câu giới thiệu 30
giây}\label{2-cuxe2u-giux1edbi-thiux1ec7u-30-giuxe2y}

\begin{quote}
Phần của em nhận telemetry từ ESP32 qua MQTT, kiểm tra và lưu vào
MongoDB Time Series. Sau khi lưu, anomaly engine áp dụng hard rule cho
gas và motion, còn nhiệt độ và độ ẩm có rolling Z-score N=30. Khi có
anomaly HIGH hoặc CRITICAL, backend có thể bật quạt qua MQTT và gửi
email, trong đó email được giới hạn một lần mỗi 15 phút theo từng cặp
thiết bị và loại cảm biến. Với CB-03, người dùng bật quạt trên React,
backend publish lệnh QoS 1, ESP32 nhận JSON và đổi GPIO 16.
\end{quote}

\section{3. Sơ đồ tổng thể phần
Kiệt}\label{3-sux1a1-ux111ux1ed3-tux1ed5ng-thux1ec3-phux1ea7n-kiux1ec7t}

\begin{center}
\resizebox{\textwidth}{!}{%
\begin{tikzpicture}[
  node distance=7mm and 6mm,
  box/.style={draw=Accent,rounded corners=2mm,fill=Panel,text=BodyText,
              align=center,font=\sffamily\bfseries\scriptsize,
              minimum height=9mm,minimum width=25mm},
  flow/.style={-{Stealth[length=2.3mm]},draw=MutedText,thick},
  label/.style={font=\sffamily\bfseries\scriptsize,text=Accent}
]
\node[label] (controltitle) {LUỒNG ĐIỀU KHIỂN};
\node[box,below=3mm of controltitle] (ui) {React\\Overview};
\node[box,right=of ui] (api) {Express\\Backend};
\node[box,right=of api] (broker) {EMQX\\Broker};
\node[box,right=of broker] (esp) {ESP32};
\node[box,right=of esp] (relay) {Relay +\\DC Fan};
\draw[flow] (ui) -- (api);
\draw[flow] (api) -- (broker);
\draw[flow] (broker) -- (esp);
\draw[flow] (esp) -- (relay);

\node[label,below=17mm of ui] (datatitle) {LUỒNG TELEMETRY VÀ PHÂN TÍCH};
\node[box,below=3mm of datatitle] (esp2) {ESP32};
\node[box,right=of esp2] (broker2) {EMQX};
\node[box,right=of broker2] (ingest) {Validate +\\ingest};
\node[box,right=of ingest] (mongo) {MongoDB\\Time Series};
\node[box,right=of mongo] (engine) {Anomaly\\Engine};
\draw[flow] (esp2) -- (broker2);
\draw[flow] (broker2) -- (ingest);
\draw[flow] (ingest) -- (mongo);
\draw[flow] (mongo) -- (engine);

\node[box,below=14mm of ingest,minimum width=31mm,draw=AccentTwo] (log) {anomaly\_logs};
\node[box,right=9mm of log,minimum width=34mm,draw=AccentTwo] (mail) {Email\\HIGH / CRITICAL};
\node[box,right=9mm of mail,minimum width=34mm,draw=AccentTwo] (mqtt) {MQTT\\fan command};
\draw[flow] (engine) -- (log);
\draw[flow] (engine) -- (mail);
\draw[flow] (engine) -- (mqtt);
\end{tikzpicture}
}
\end{center}

Ba ý cần nhấn mạnh:

\begin{enumerate}
\def\labelenumi{\arabic{enumi}.}
\tightlist
\item
  MQTT tách backend khỏi ESP32, vì vậy web không điều khiển GPIO trực
  tiếp.
\item
  Telemetry được validate trước khi lưu và trước khi chạy anomaly
  engine.
\item
  Gas vẫn có cảnh báo cục bộ trên ESP32; backend không được xem là lớp
  an toàn duy nhất.
\end{enumerate}

\section{4. CB-03 - Điều khiển quạt từ
Website}\label{4-cb-03---ux111iux1ec1u-khiux1ec3n-quux1ea1t-tux1eeb-website}

\subsection{4.1 Luồng xử lý}\label{41-luux1ed3ng-xux1eed-luxfd}

\begin{enumerate}
\def\labelenumi{\arabic{enumi}.}
\tightlist
\item
  Người dùng nhấn nút fan trên \texttt{Overview.jsx}.
\item
  Frontend đảo trạng thái hiện tại và gửi
  \nolinkurl{POST /api/v1/control/actuator}.
\item
  JWT được gắn vào request bởi \texttt{apiFetch()}.
\item
  Backend kiểm tra \texttt{deviceId}, actuator và state.
\item
  \texttt{mqttService.publishControlCommand()} publish JSON lên
  \nolinkurl{safebox/{deviceId}/control}, QoS 1.
\item
  \texttt{MqttClient} trên ESP32 parse payload rồi gọi callback
  \texttt{onControlMessage()} trong \texttt{main.cpp}.
\item
  Callback định tuyến actuator \texttt{fan} sang module \texttt{Relay};
  actuator lạ bị bỏ qua.
\item
  \texttt{Relay::on()} hoặc \texttt{Relay::off()} ghi GPIO 16 theo cấu
  hình \texttt{RELAY\_ACTIVE\_LOW}.
\item
  Backend tạm ghi trạng thái yêu cầu để UI phản hồi nhanh; telemetry kế
  tiếp từ ESP32 xác nhận lại trạng thái thật.
\end{enumerate}

\subsection{4.2 Contract quan trọng}\label{42-contract-quan-trux1ecdng}

Request từ frontend:

\begin{Shaded}
\begin{Highlighting}[]
\FunctionTok{\{}
  \DataTypeTok{"deviceId"}\FunctionTok{:} \StringTok{"safebox{-}01"}\FunctionTok{,}
  \DataTypeTok{"actuator"}\FunctionTok{:} \StringTok{"fan"}\FunctionTok{,}
  \DataTypeTok{"state"}\FunctionTok{:} \KeywordTok{true}
\FunctionTok{\}}
\end{Highlighting}
\end{Shaded}

MQTT command gửi xuống ESP32:

\begin{Shaded}
\begin{Highlighting}[]
\FunctionTok{\{}
  \DataTypeTok{"actuator"}\FunctionTok{:} \StringTok{"fan"}\FunctionTok{,}
  \DataTypeTok{"state"}\FunctionTok{:} \StringTok{"on"}
\FunctionTok{\}}
\end{Highlighting}
\end{Shaded}

Topic:

\begin{Shaded}
\begin{Highlighting}[]
\NormalTok{safebox/safebox{-}01/control}
\end{Highlighting}
\end{Shaded}

\subsection{4.3 Vì sao REST dùng boolean nhưng MQTT dùng
on/off?}\label{43-vuxec-sao-rest-duxf9ng-boolean-nhux1b0ng-mqtt-duxf9ng-onoff}

\begin{itemize}
\tightlist
\item
  Boolean thuận tiện cho toggle React và validation API.
\item
  Chuỗi \texttt{on}/\texttt{off} dễ đọc trên MQTT client và Serial
  Monitor.
\item
  Backend là nơi chuyển đổi contract, nên ESP32 không phải hiểu nhiều
  định dạng.
\end{itemize}

\subsection{4.4 Vì sao dùng QoS 1?}\label{44-vuxec-sao-duxf9ng-qos-1}

QoS 1 bảo đảm lệnh được giao ít nhất một lần giữa client và broker. Lệnh
bật/tắt là idempotent: nhận \texttt{fan=on} hai lần vẫn chỉ cho fan ở
trạng thái ON. Không được nói QoS 1 là ``exactly once''; đó là QoS 2.

\subsection{4.5 Requested state và confirmed
state}\label{45-requested-state-vuxe0-confirmed-state}

Sau khi broker nhận publish thành công, backend cập nhật requested state
để API status phản hồi ngay. Tuy nhiên đây chưa phải bằng chứng relay đã
đổi trạng thái. Telemetry kế tiếp chứa \texttt{actuatorStates.fanOn} sẽ
reconcile lại state theo báo cáo của ESP32.

Đây là câu trả lời tốt khi bị hỏi ``publish thành công có chắc quạt quay
chưa?'':

\begin{quote}
Chưa chắc. Publish callback chỉ chứng minh broker nhận lệnh. Xác nhận
thiết bị cần telemetry phản hồi hoặc acknowledgment riêng. Bản hiện tại
dùng telemetry kế tiếp để reconcile; nếu cần xác nhận mạnh hơn có thể
thêm command ID và ACK topic.
\end{quote}

\subsection{4.6 Fail-safe của
firmware}\label{46-fail-safe-cux1ee7a-firmware}

\begin{itemize}
\tightlist
\item
  GPIO 16 được cấu hình OUTPUT khi khởi động.
\item
  \texttt{Relay::begin()} gọi \texttt{off()} nên fan mặc định OFF; với
  HL-52S active-low, mức boot là HIGH.
\item
  JSON sai bị bỏ qua; actuator lạ không có handler. Backend chỉ publish
  state hợp lệ \texttt{on}/\texttt{off}.
\item
  Không có local auto-trigger cho fan ngoài yêu cầu backend, đúng phạm
  vi dự án.
\end{itemize}

\section{5. NC-01 - Anomaly engine và
Z-score}\label{5-nc-01---anomaly-engine-vuxe0-z-score}

\subsection{5.1 Công thức}\label{51-cuxf4ng-thux1ee9c}

\begin{Shaded}
\begin{Highlighting}[]
\NormalTok{mean = tổng 30 giá trị / 30}
\NormalTok{variance = tổng (x\_i {-} mean)\^{}2 / 30}
\NormalTok{standard deviation = sqrt(variance)}
\NormalTok{Z = (giá trị mới {-} mean) / standard deviation}
\end{Highlighting}
\end{Shaded}

Một điểm được xem là anomaly thống kê khi đồng thời:

\begin{Shaded}
\begin{Highlighting}[]
\NormalTok{|Z| \textgreater{} zScoreLimit}
\NormalTok{và}
\NormalTok{|giá trị mới {-} mean| \textgreater{} minimum physical delta}
\end{Highlighting}
\end{Shaded}

Giá trị mặc định:

{\def\LTcaptype{none} % do not increment counter
\begin{longtable}[]{@{}
  >{\raggedright\arraybackslash}p{(\linewidth - 2\tabcolsep) * \real{0.3100}}
  >{\raggedleft\arraybackslash}p{(\linewidth - 2\tabcolsep) * \real{0.6400}}@{}}
\toprule\noalign{}
\begin{minipage}[b]{\linewidth}\raggedright
Thông số
\end{minipage} & \begin{minipage}[b]{\linewidth}\raggedleft
Giá trị
\end{minipage} \\
\midrule\noalign{}
\endhead
\bottomrule\noalign{}
\endlastfoot
Window & 30 mẫu trước đó \\
Z-score limit & 3.0 \\
Minimum delta nhiệt độ & 1.0°C \\
Minimum delta độ ẩm & 2.0\% \\
\end{longtable}
}

\subsection{5.2 Vì sao dùng 30 mẫu trước, không cho mẫu mới vào
trước?}\label{52-vuxec-sao-duxf9ng-30-mux1eabu-trux1b0ux1edbc-khuxf4ng-cho-mux1eabu-mux1edbi-vuxe0o-trux1b0ux1edbc}

Mẫu mới phải được so với baseline trước đó. Nếu đưa spike vào window
trước khi tính mean và độ lệch chuẩn, chính spike sẽ kéo baseline về
phía nó và làm giảm Z-score. Code tính thống kê trước, sau đó mới
\texttt{push} mẫu mới và bỏ mẫu cũ nhất.

\subsection{5.3 Cold start và zero
variance}\label{53-cold-start-vuxe0-zero-variance}

\begin{itemize}
\tightlist
\item
  Ít hơn 30 mẫu: chưa đủ baseline, trả \texttt{ready=false} và dựa vào
  hard threshold.
\item
  Độ lệch chuẩn bằng 0: không chia cho 0; trả Z=0 và không tạo anomaly
  thống kê.
\item
  Hard threshold vẫn luôn hoạt động, kể cả khi Z-score chưa sẵn sàng.
\end{itemize}

\subsection{5.4 Vì sao gas và motion không dùng
Z-score?}\label{54-vuxec-sao-gas-vuxe0-motion-khuxf4ng-duxf9ng-z-score}

\begin{itemize}
\tightlist
\item
  Nhiệt độ và độ ẩm là biến liên tục, phù hợp với phân tích độ lệch
  quanh baseline.
\item
  Motion là boolean, không có ý nghĩa Gaussian.
\item
  Gas là biến an toàn có ngưỡng cứng 600 PPM; chờ đủ 30 mẫu hoặc chờ
  Z-score có thể gây nguy hiểm.
\end{itemize}

\subsection{5.5 Rule matrix thực
tế}\label{55-rule-matrix-thux1ef1c-tux1ebf}

{\def\LTcaptype{none} % do not increment counter
\begin{longtable}[]{@{}
  >{\raggedright\arraybackslash}p{(\linewidth - 8\tabcolsep) * \real{0.1400}}
  >{\raggedright\arraybackslash}p{(\linewidth - 8\tabcolsep) * \real{0.2200}}
  >{\raggedright\arraybackslash}p{(\linewidth - 8\tabcolsep) * \real{0.1800}}
  >{\raggedright\arraybackslash}p{(\linewidth - 8\tabcolsep) * \real{0.1400}}
  >{\raggedright\arraybackslash}p{(\linewidth - 8\tabcolsep) * \real{0.2700}}@{}}
\toprule\noalign{}
\begin{minipage}[b]{\linewidth}\raggedright
Sensor
\end{minipage} & \begin{minipage}[b]{\linewidth}\raggedright
Điều kiện
\end{minipage} & \begin{minipage}[b]{\linewidth}\raggedright
Trigger
\end{minipage} & \begin{minipage}[b]{\linewidth}\raggedright
Severity
\end{minipage} & \begin{minipage}[b]{\linewidth}\raggedright
Hành động
\end{minipage} \\
\midrule\noalign{}
\endhead
\bottomrule\noalign{}
\endlastfoot
Gas & \texttt{gasLevel\ \textgreater{}\ gasMax} & Hard threshold &
CRITICAL & Log + email \\
Temperature & Ngoài min/max & Hard threshold & HIGH & Log + email; bật
fan nếu vượt max \\
Temperature & Auto mode và Z bất thường & Z-score & HIGH & Log + email;
bật fan nếu spike hướng lên \\
Motion & Motion=true và system armed & Hard threshold & HIGH & Log +
email + alert LED \\
Humidity & Ngoài min/max & Hard threshold & MEDIUM & Log \\
Humidity & Auto mode và Z bất thường & Z-score & MEDIUM & Log \\
\end{longtable}
}

\subsection{5.6 Dedupe khác email throttle thế
nào?}\label{56-dedupe-khuxe1c-email-throttle-thux1ebf-nuxe0o}

\begin{itemize}
\tightlist
\item
  Anomaly dedupe mặc định 60 giây, key gồm
  \texttt{deviceId:sensorType:triggerType}. Mục đích là tránh ghi cùng
  sự kiện quá dày.
\item
  Email throttle là 15 phút, key gồm \texttt{deviceId:sensorType}. Mục
  đích là tránh spam hộp thư.
\item
  Hai cơ chế giải quyết hai vấn đề khác nhau và không nên gộp.
\end{itemize}

\subsection{5.7 Thứ tự xử lý anomaly có cố định
không?}\label{57-thux1ee9-tux1ef1-xux1eed-luxfd-anomaly-cuxf3-cux1ed1-ux111ux1ecbnh-khuxf4ng}

Không. Các candidate được xử lý bằng \texttt{Promise.all}, nên email
hoặc MQTT của từng candidate có thể hoàn tất khác thứ tự. Điều cần bảo
đảm là đủ tập anomaly và đúng nội dung, không dựa vào thứ tự hoàn thành.

\section{6. NC-03 - MongoDB Time
Series}\label{6-nc-03---mongodb-time-series}

\subsection{6.1 Tại sao dùng Time Series
collection?}\label{61-tux1ea1i-sao-duxf9ng-time-series-collection}

Telemetry là dữ liệu liên tục, mỗi document có timestamp và deviceId.
Time Series tối ưu cách MongoDB lưu và truy vấn dữ liệu theo thời gian,
đồng thời hỗ trợ retention tự động.

\subsection{6.2 Cấu hình
authoritative}\label{62-cux1ea5u-huxecnh-authoritative}

\begin{Shaded}
\begin{Highlighting}[]
\NormalTok{collection: sensor\_readings}
\NormalTok{timeField: timestamp}
\NormalTok{metaField: deviceId}
\NormalTok{granularity: seconds}
\NormalTok{expireAfterSeconds: 2,592,000 (30 ngày)}
\end{Highlighting}
\end{Shaded}

\subsection{6.3 Luồng ghi dữ
liệu}\label{63-luux1ed3ng-ghi-dux1eef-liux1ec7u}

\begin{enumerate}
\def\labelenumi{\arabic{enumi}.}
\tightlist
\item
  \texttt{mqttService} nhận topic telemetry.
\item
  \texttt{validateTelemetryPayload()} kiểm tra device ID, timestamp,
  metric và actuator state.
\item
  \texttt{SensorReading.create()} ghi document vào
  \texttt{sensor\_readings}.
\item
  Document vừa ghi được chuyển cho anomaly engine.
\end{enumerate}

Việc lưu trước rồi phân tích giúp giữ lại dữ liệu gốc ngay cả khi một
hành động anomaly như email hoặc MQTT gặp lỗi.

\subsection{\texorpdfstring{6.4 Vì sao schema đặt
\texttt{autoCreate=false}?}{6.4 Vì sao schema đặt autoCreate=false?}}\label{64-vuxec-sao-schema-ux111ux1eb7t-autocreatefalse}

Mongoose không tự tạo collection thường. \texttt{timeSeriesService} chủ
động tạo collection với đúng \texttt{timeseries} và TTL option trước khi
model ghi dữ liệu. Nếu collection đã tồn tại nhưng sai cấu hình, backend
fail fast và yêu cầu migration thay vì âm thầm chạy sai.

\subsection{6.5 Historical query lấy dữ liệu thế
nào?}\label{65-historical-query-lux1ea5y-dux1eef-liux1ec7u-thux1ebf-nuxe0o}

Backend lọc theo device và khoảng thời gian, sort giảm dần để lấy đúng N
mẫu mới nhất, áp dụng limit, sau đó reverse trong bộ nhớ để frontend
nhận chuỗi theo thời gian tăng dần.

\begin{Shaded}
\begin{Highlighting}[]
\NormalTok{MongoDB: newest {-}\textgreater{} oldest, limit N}
\NormalTok{Response: oldest {-}\textgreater{} newest trong N mẫu vừa chọn}
\end{Highlighting}
\end{Shaded}

Nếu sort tăng dần trước limit, API sẽ trả N mẫu cũ nhất trong khoảng,
không phải N mẫu mới nhất.

\subsection{6.6 TTL có xóa đúng ngay giây thứ 2,592,000
không?}\label{66-ttl-cuxf3-xuxf3a-ux111uxfang-ngay-giuxe2y-thux1ee9-2592000-khuxf4ng}

Không nên khẳng định chính xác tới từng giây. MongoDB TTL monitor chạy
nền, nên document được xóa sau khi hết hạn với một độ trễ nhỏ. Ý nghĩa
là retention tự động, không phải timer realtime.

\section{7. NC-05 - Email
Notification}\label{7-nc-05---email-notification}

\subsection{7.1 Luồng gửi email}\label{71-luux1ed3ng-gux1eedi-email}

\begin{enumerate}
\def\labelenumi{\arabic{enumi}.}
\tightlist
\item
  Anomaly candidate được tạo.
\item
  \texttt{executeActions()} chỉ gọi email cho severity HIGH hoặc
  CRITICAL.
\item
  \texttt{sendAlertEmail()} đọc cấu hình SMTP sender từ environment.
\item
  Recipient resolver lấy \texttt{User.email} của các tài khoản có
  \texttt{settings.emailEnabled=true}.
\item
  Hàm tạo key \texttt{deviceId:sensorType} và kiểm tra lần gửi gần nhất.
\item
  Nếu chưa bị throttle, Nodemailer gửi cả text và HTML email.
\item
  Chỉ sau khi \texttt{sendMail()} thành công mới lưu thời điểm gửi.
\end{enumerate}

\subsection{7.2 Vì sao chỉ HIGH và
CRITICAL?}\label{72-vuxec-sao-chux1ec9-high-vuxe0-critical}

MEDIUM như humidity bất thường vẫn được ghi log và hiển thị trên
dashboard nhưng không gửi email để tránh cảnh báo quá mức. HIGH và
CRITICAL cần người dùng phản ứng sớm hơn.

\subsection{7.3 Vì sao throttle theo cặp device và
sensor?}\label{73-vuxec-sao-throttle-theo-cux1eb7p-device-vuxe0-sensor}

Gas trên \texttt{safebox-01} không được chặn email nhiệt độ trên cùng
thiết bị, và cũng không được chặn gas trên \texttt{safebox-02}. Vì vậy
key phải chứa cả device và sensor.

Ví dụ:

\begin{Shaded}
\begin{Highlighting}[]
\NormalTok{safebox{-}01:gasLevel       {-}\textgreater{} một window riêng}
\NormalTok{safebox{-}01:temperature    {-}\textgreater{} một window riêng}
\NormalTok{safebox{-}02:gasLevel       {-}\textgreater{} một window riêng}
\end{Highlighting}
\end{Shaded}

\subsection{7.4 Vì sao lưu throttle sau khi gửi thành
công?}\label{74-vuxec-sao-lux1b0u-throttle-sau-khi-gux1eedi-thuxe0nh-cuxf4ng}

Nếu SMTP lỗi mà vẫn ghi thời điểm throttle, các lần retry tiếp theo sẽ
bị chặn dù người dùng chưa nhận email nào. Code chỉ cập nhật map sau
\texttt{sendMail()} thành công.

\subsection{7.5 Giới hạn cần nói
thật}\label{75-giux1edbi-hux1ea1n-cux1ea7n-nuxf3i-thux1eadt}

\begin{itemize}
\tightlist
\item
  Throttle map nằm trong RAM, nên reset khi backend restart.
\item
  Recipient lấy từ email đăng ký của mọi user đang bật email trong
  Settings; nếu tất cả user tắt thì SMTP không được gọi.
\item
  Hệ thống chưa có quan hệ device-owner, nên mọi user đang bật đều nhận
  anomaly của thiết bị. Bản production cần ánh xạ user với device.
\item
  \texttt{ALERT\_RECIPIENT\_EMAIL} chỉ là fallback khi database chưa có
  tài khoản nào.
\item
  Nếu triển khai nhiều instance, nên chuyển throttle sang Redis/MongoDB.
\end{itemize}

\section{8. Bản đồ file cần mở khi vấn
đáp}\label{8-bux1ea3n-ux111ux1ed3-file-cux1ea7n-mux1edf-khi-vux1ea5n-ux111uxe1p}

{\def\LTcaptype{none} % do not increment counter
\begin{longtable}[]{@{}
  >{\raggedright\arraybackslash}p{(\linewidth - 4\tabcolsep) * \real{0.2300}}
  >{\raggedright\arraybackslash}p{(\linewidth - 4\tabcolsep) * \real{0.3700}}
  >{\raggedright\arraybackslash}p{(\linewidth - 4\tabcolsep) * \real{0.3500}}@{}}
\toprule\noalign{}
\begin{minipage}[b]{\linewidth}\raggedright
Phần
\end{minipage} & \begin{minipage}[b]{\linewidth}\raggedright
File chính
\end{minipage} & \begin{minipage}[b]{\linewidth}\raggedright
Vai trò
\end{minipage} \\
\midrule\noalign{}
\endhead
\bottomrule\noalign{}
\endlastfoot
CB-03 frontend & \nolinkurl{frontend/src/pages/Overview.jsx} & Toggle,
POST API, optimistic state \\
CB-03 auth helper & \nolinkurl{frontend/src/api.js} & Gắn JWT vào
request \\
CB-03 API & \nolinkurl{backend/routes/control.js} & Validate request và
gọi MQTT service \\
CB-03 state & \nolinkurl{backend/services/actuatorStateService.js} &
Requested state và telemetry reconciliation \\
CB-03 MQTT & \nolinkurl{backend/services/mqttService.js} & Publish QoS 1
và nhận telemetry/status \\
CB-03 firmware & \nolinkurl{firmware/src/main.cpp} & Đăng ký callback và
định tuyến actuator \\
CB-03 MQTT parser & \nolinkurl{firmware/src/mqtt/mqtt_client.cpp} &
Subscribe control topic và parse JSON \\
CB-03 relay & \nolinkurl{firmware/src/actuators/relay.cpp} & Ghi GPIO 16
theo polarity cấu hình \\
NC-01 & \nolinkurl{backend/services/anomalyEngine.js} & Z-score, rules,
dedupe, actions \\
NC-01 schema & \nolinkurl{backend/models/Schemas.js} & Cấu trúc
\texttt{anomaly\_logs} \\
NC-03 & \nolinkurl{backend/services/timeSeriesService.js} & Tạo và
verify Time Series collection \\
NC-03 ingestion & \nolinkurl{backend/services/telemetryService.js} &
Validate, create, gọi anomaly engine \\
NC-03 API & \nolinkurl{backend/routes/telemetry.js} & Latest và
historical query \\
NC-05 & \nolinkurl{backend/services/emailService.js} & SMTP, template,
severity, throttle \\
Khởi động & \nolinkurl{backend/server.js} & MongoDB, Time Series, email,
MQTT, mock \\
\end{longtable}
}

Chi tiết code có trong \url{THAI_KIET_CODE_WALKTHROUGH.md}.

\section{9. Bằng chứng hiện
có}\label{9-bux1eb1ng-chux1ee9ng-hiux1ec7n-cuxf3}

Chạy tại thư mục \texttt{backend}:

\begin{Shaded}
\begin{Highlighting}[]
\NormalTok{npm run test}\OperatorTok{:}\NormalTok{kiet}
\end{Highlighting}
\end{Shaded}

Bộ test chứng minh:

\begin{itemize}
\tightlist
\item
  Mapping và reconcile actuator state.
\item
  JWT chặn request không xác thực trong CB-03 simulation.
\item
  N=30, per-device buffer, zero variance và minimum delta.
\item
  Gas/motion không đi qua Z-score.
\item
  Hard-rule matrix và hành động alert LED.
\item
  Temperature spike tạo Z-score anomaly, bật fan, ghi log và gửi email.
\item
  Email throttle cùng key và độc lập giữa các key.
\item
  Time Series options và retention 30 ngày.
\item
  Historical API chọn N mẫu mới nhất và trả theo thứ tự thời gian.
\end{itemize}

Điều chưa được test phần mềm thay thế hoàn toàn:

\begin{itemize}
\tightlist
\item
  Relay vật lý có click và fan quay dưới 500 ms hay không.
\item
  SMTP thật có giao email trong dưới 5 giây hay không.
\item
  EMQX và Atlas thật có ổn định trong phiên demo hay không.
\end{itemize}

Các mục này phải được đo bằng TC-02, TC-04, TC-05 và TC-06 với hạ tầng
thật.

\section{10. Cheat sheet một phút trước khi vào
phòng}\label{10-cheat-sheet-mux1ed9t-phuxfat-trux1b0ux1edbc-khi-vuxe0o-phuxf2ng}

\begin{Shaded}
\begin{Highlighting}[]
\NormalTok{CB{-}03:}
\NormalTok{React POST boolean {-}\textgreater{} backend validate {-}\textgreater{} MQTT on/off QoS 1}
\NormalTok{{-}\textgreater{} ESP32 MqttClient {-}\textgreater{} onControlMessage {-}\textgreater{} Relay {-}\textgreater{} GPIO 16}
\NormalTok{Publish success != physical ACK; telemetry kế tiếp reconcile state.}

\NormalTok{NC{-}01:}
\NormalTok{Z{-}score chỉ temp/humidity, previous N=30, |Z|\textgreater{}3 và min delta.}
\NormalTok{Gas \textgreater{}600 và motion dùng hard rule.}
\NormalTok{Cold start dùng hard bounds; sigma=0 thì Z=0.}
\NormalTok{Dedupe anomaly 60s khác email throttle 15m.}

\NormalTok{NC{-}03:}
\NormalTok{sensor\_readings Time Series; timestamp=timeField; deviceId=metaField.}
\NormalTok{Granularity seconds; TTL 30 ngày.}
\NormalTok{Validate {-}\textgreater{} create {-}\textgreater{} anomaly engine.}
\NormalTok{Sort DESC + limit + reverse để lấy N mẫu mới nhất theo thứ tự tăng dần.}

\NormalTok{NC{-}05:}
\NormalTok{Nodemailer chỉ HIGH/CRITICAL.}
\NormalTok{Throttle key=deviceId:sensorType, 15 phút.}
\NormalTok{Chỉ ghi lastSent sau khi SMTP gửi thành công.}

\NormalTok{MQ2 telemetry:}
\NormalTok{48h burn{-}in {-}\textgreater{} stable clean{-}air Ro {-}\textgreater{} NVS {-}\textgreater{} 5m warm{-}up on later boots.}
\NormalTok{Only ready has numeric gasLevel; preparing sends null + sensorStatus.}
\NormalTok{100ms high{-}priority local task; raw fallback stays local, no fake email.}
\end{Highlighting}
\end{Shaded}

\chapter{Bộ câu hỏi vấn đáp phần Thái
Kiệt}\label{bux1ed9-cuxe2u-hux1ecfi-vux1ea5n-ux111uxe1p-phux1ea7n-thuxe1i-kiux1ec7t}

Câu trả lời mẫu được viết theo cách nói trực tiếp. Không cần học thuộc
từng chữ; cần hiểu luận điểm và chỉ được code minh họa.

\section{A. Câu hỏi tổng quan}\label{a-cuxe2u-hux1ecfi-tux1ed5ng-quan}

\subsection{1. Em phụ trách phần
nào?}\label{1-em-phux1ee5-truxe1ch-phux1ea7n-nuxe0o}

Em phụ trách CB-03 điều khiển quạt từ website, NC-01 anomaly detection
bằng Z-score và hard rules, NC-03 lưu telemetry bằng MongoDB Time
Series, và NC-05 gửi email cảnh báo có throttle.

\subsection{2. Bốn phần này liên hệ với nhau thế
nào?}\label{2-bux1ed1n-phux1ea7n-nuxe0y-liuxean-hux1ec7-vux1edbi-nhau-thux1ebf-nuxe0o}

NC-03 nhận và lưu telemetry. NC-01 phân tích chính telemetry đó. NC-01
có thể gọi lại luồng CB-03 để bật fan khi nhiệt độ cao và gọi NC-05 để
gửi email khi severity HIGH hoặc CRITICAL.

\subsection{3. Thành phần nào là trung tâm điều
phối?}\label{3-thuxe0nh-phux1ea7n-nuxe0o-luxe0-trung-tuxe2m-ux111iux1ec1u-phux1ed1i}

Express backend. Nó nhận REST từ frontend, giao tiếp MQTT với ESP32, ghi
MongoDB, chạy anomaly engine và gọi Nodemailer.

\subsection{4. Vì sao không cho frontend kết nối MQTT trực
tiếp?}\label{4-vuxec-sao-khuxf4ng-cho-frontend-kux1ebft-nux1ed1i-mqtt-trux1ef1c-tiux1ebfp}

Backend cần giữ MQTT credential, validate quyền người dùng, chuẩn hóa
lệnh và tập trung business rule. Nếu frontend kết nối broker trực tiếp
thì credential bị lộ và khó kiểm soát quyền publish.

\subsection{5. Điểm an toàn quan trọng nhất của kiến trúc là
gì?}\label{5-ux111iux1ec3m-an-touxe0n-quan-trux1ecdng-nhux1ea5t-cux1ee7a-kiux1ebfn-truxfac-luxe0-guxec}

Cảnh báo gas nguy hiểm vẫn phải chạy cục bộ trên ESP32 dưới 200 ms,
không phụ thuộc backend. Backend bổ sung lưu log và email chứ không thay
thế edge safety.

\subsection{6. Em chứng minh code chạy bằng
gì?}\label{6-em-chux1ee9ng-minh-code-chux1ea1y-bux1eb1ng-guxec}

Em dùng test tự động cho logic và integration simulation, cộng với
firmware/frontend build. Riêng latency relay, SMTP thật, EMQX và Atlas
phải đo bằng test end-to-end với hạ tầng thật.

\subsection{7. Nếu được làm lại, em ưu tiên cải tiến
gì?}\label{7-nux1ebfu-ux111ux1b0ux1ee3c-luxe0m-lux1ea1i-em-ux1b0u-tiuxean-cux1ea3i-tiux1ebfn-guxec}

Em sẽ thêm command ID và ACK topic cho actuator, chuyển dedupe/throttle
state sang Redis hoặc MongoDB để không mất khi restart, và gắn device
ownership với user cho mô hình nhiều thiết bị.

\section{B. CB-03 - Điều khiển
quạt}\label{b-cb-03---ux111iux1ec1u-khiux1ec3n-quux1ea1t}

\subsection{8. Luồng bật quạt bắt đầu từ
đâu?}\label{8-luux1ed3ng-bux1eadt-quux1ea1t-bux1eaft-ux111ux1ea7u-tux1eeb-ux111uxe2u}

Từ
\texttt{handleActuatorToggle(\textquotesingle{}fan\textquotesingle{},\ actuators.fanOn)}
trong \texttt{Overview.jsx}. Frontend gửi POST tới
\nolinkurl{/api/v1/control/actuator} với deviceId, actuator và boolean
state.

\subsection{9. Request có được xác thực
không?}\label{9-request-cuxf3-ux111ux1b0ux1ee3c-xuxe1c-thux1ef1c-khuxf4ng}

Có. \texttt{apiFetch()} gắn JWT vào Authorization header. Backend mount
\texttt{authenticateToken} trước control router. Không có token trả 401,
token sai hoặc hết hạn trả 403.

\subsection{10. Backend validate những
gì?}\label{10-backend-validate-nhux1eefng-guxec}

Route yêu cầu actuator và state, actuator phải nằm trong allowlist,
state phải là boolean hoặc chuỗi on/off. MQTT service kiểm tra thêm
deviceId, actuator và boolean trước khi publish.

\subsection{11. Topic điều khiển là
gì?}\label{11-topic-ux111iux1ec1u-khiux1ec3n-luxe0-guxec}

\nolinkurl{safebox/{deviceId}/control}, ví dụ
\nolinkurl{safebox/safebox-01/control}.

\subsection{12. Payload MQTT là gì?}\label{12-payload-mqtt-luxe0-guxec}

\texttt{\{"actuator":"fan","state":"on"\}} hoặc state \texttt{off}.

\subsection{13. Vì sao dùng QoS 1?}\label{13-vuxec-sao-duxf9ng-qos-1}

Lệnh điều khiển quan trọng hơn telemetry thông thường, nên cần
at-least-once delivery. Bật/tắt là idempotent nên duplicate không gây
toggle ngược.

\subsection{14. QoS 1 có bảo đảm quạt đã quay
không?}\label{14-qos-1-cuxf3-bux1ea3o-ux111ux1ea3m-quux1ea1t-ux111uxe3-quay-khuxf4ng}

Không. Nó chỉ bảo đảm giao message ở tầng MQTT đến broker/client theo cơ
chế QoS. Muốn biết phần cứng đã thực thi phải dùng telemetry phản hồi
hoặc ACK riêng từ ESP32.

\subsection{15. Vì sao control message không
retained?}\label{15-vuxec-sao-control-message-khuxf4ng-retained}

Nếu retained, ESP32 kết nối lại có thể nhận lệnh cũ và bật quạt ngoài ý
muốn. Control command là event hiện tại nên dùng \texttt{retain=false}.

\subsection{16. Nếu MQTT offline thì API trả
gì?}\label{16-nux1ebfu-mqtt-offline-thuxec-api-trux1ea3-guxec}

\texttt{publishControlCommand()} ném lỗi có code \texttt{MQTT\_OFFLINE},
route chuyển thành HTTP 503. Frontend hiển thị control feedback lỗi và
không giả vờ thành công.

\subsection{17. Trạng thái fan trên backend lấy từ
đâu?}\label{17-trux1ea1ng-thuxe1i-fan-truxean-backend-lux1ea5y-tux1eeb-ux111uxe2u}

Ngay sau publish thành công, backend lưu requested state trong
\texttt{actuatorStateService}. Khi telemetry ESP32 về, snapshot được
synchronize lại bằng \texttt{actuatorStates} do thiết bị báo cáo.

\subsection{\texorpdfstring{18. Vì sao tách
\texttt{actuatorStateService} khỏi mock
generator?}{18. Vì sao tách actuatorStateService khỏi mock generator?}}\label{18-vuxec-sao-tuxe1ch-actuatorstateservice-khux1ecfi-mock-generator}

State này được dùng cả simulation và MQTT thật. Đặt nó trong file mock
gây hiểu nhầm ownership. Service riêng làm rõ đây là trạng thái runtime
chung; mock generator chỉ sinh dữ liệu.

\subsection{19. Firmware chống payload sai thế
nào?}\label{19-firmware-chux1ed1ng-payload-sai-thux1ebf-nuxe0o}

ArduinoJson parse theo buffer và length; JSON lỗi bị bỏ qua.
\texttt{main.cpp} chỉ định tuyến \texttt{fan}, \texttt{buzzer},
\texttt{alertLed}, nên actuator lạ không tạo tác động. Backend đã
validate và chỉ publish state hợp lệ; phía firmware xử lý chuỗi khác
\texttt{on} theo hướng fail-off.

\subsection{20. Fan ở trạng thái nào khi ESP32
boot?}\label{20-fan-ux1edf-trux1ea1ng-thuxe1i-nuxe0o-khi-esp32-boot}

OFF. \texttt{Relay::begin()} cấu hình pin OUTPUT rồi gọi \texttt{off()}.
Với HL-52S mặc định active-low, OFF tương ứng GPIO 16 HIGH; polarity có
thể đổi bằng \texttt{RELAY\_ACTIVE\_LOW}.

\subsection{21. Nếu QoS 1 giao cùng lệnh hai lần thì
sao?}\label{21-nux1ebfu-qos-1-giao-cuxf9ng-lux1ec7nh-hai-lux1ea7n-thuxec-sao}

Không sao vì command mang trạng thái đích, không phải lệnh toggle. Gọi
\texttt{Relay::on()} hai lần vẫn cho fan ON.

\subsection{22. Hệ thống có hỗ trợ nhiều device hoàn chỉnh
chưa?}\label{22-hux1ec7-thux1ed1ng-cuxf3-hux1ed7-trux1ee3-nhiux1ec1u-device-houxe0n-chux1ec9nh-chux1b0a}

Topic, telemetry schema và historical query có deviceId. Tuy nhiên
runtime actuator snapshot hiện là một object vì coursework dùng một
device chính. Muốn multi-device hoàn chỉnh cần đổi thành
\texttt{Map\textless{}deviceId,\ actuatorStates\textgreater{}}.

\section{C. NC-01 - Z-score và
anomaly}\label{c-nc-01---z-score-vuxe0-anomaly}

\subsection{23. Công thức Z-score là
gì?}\label{23-cuxf4ng-thux1ee9c-z-score-luxe0-guxec}

\nolinkurl{Z = (x - mean) / standardDeviation}. Code lấy mean và
standard deviation từ 30 mẫu trước đó của cùng device và cùng sensor.

\subsection{24. Vì sao chọn N=30?}\label{24-vuxec-sao-chux1ecdn-n30}

Đây là kích thước rolling window theo thiết kế dự án: đủ tạo baseline
ngắn hạn nhưng vẫn phản ứng với thay đổi môi trường. Với chu kỳ 2 giây,
30 mẫu đại diện khoảng một phút.

\subsection{25. Tại sao dùng 30 mẫu trước thay vì gồm cả mẫu
mới?}\label{25-tux1ea1i-sao-duxf9ng-30-mux1eabu-trux1b0ux1edbc-thay-vuxec-gux1ed3m-cux1ea3-mux1eabu-mux1edbi}

Mẫu mới là đối tượng cần kiểm tra. Nếu đưa nó vào baseline trước, spike
làm mean và variance tăng, khiến Z-score bị giảm và khó phát hiện hơn.

\subsection{26. Điều kiện anomaly thống kê là
gì?}\label{26-ux111iux1ec1u-kiux1ec7n-anomaly-thux1ed1ng-kuxea-luxe0-guxec}

Hai điều kiện đồng thời: trị tuyệt đối Z lớn hơn threshold mặc định 3.0
và độ lệch vật lý lớn hơn minimum delta, 1°C cho nhiệt độ hoặc 2\% cho
độ ẩm.

\subsection{27. Minimum delta để làm
gì?}\label{27-minimum-delta-ux111ux1ec3-luxe0m-guxec}

Khi standard deviation rất nhỏ, một dao động rất nhỏ cũng có thể tạo
Z-score lớn. Minimum delta lọc các spike toán học nhưng không có ý nghĩa
vật lý.

\subsection{28. Z-score áp dụng cho sensor
nào?}\label{28-z-score-uxe1p-dux1ee5ng-cho-sensor-nuxe0o}

Chỉ temperature và humidity. \texttt{MIN\_DELTA} cũng đóng vai trò
allowlist; sensor khác trả \texttt{ready=false} ngay.

\subsection{29. Vì sao gas không dùng
Z-score?}\label{29-vuxec-sao-gas-khuxf4ng-duxf9ng-z-score}

Gas là thông số an toàn có ngưỡng cứng. Gas vượt 600 PPM phải cảnh báo
ngay, không thể chờ 30 mẫu hoặc phụ thuộc phân phối thống kê.

\subsection{30. Vì sao motion không dùng
Z-score?}\label{30-vuxec-sao-motion-khuxf4ng-duxf9ng-z-score}

Motion là boolean, không phải biến liên tục có phân phối phù hợp để tính
mean và standard deviation theo cách này.

\subsection{31. Cold start xử lý thế
nào?}\label{31-cold-start-xux1eed-luxfd-thux1ebf-nuxe0o}

Khi buffer dưới 30 mẫu, Z-score chưa ready. Hard bounds vẫn chạy để
không bỏ qua nhiệt độ, độ ẩm hoặc gas nguy hiểm.

\subsection{32. Standard deviation bằng 0 thì
sao?}\label{32-standard-deviation-bux1eb1ng-0-thuxec-sao}

Phép chia không xác định nên code trả Z=0 và không tạo statistical
anomaly. Hard thresholds vẫn là lớp fallback. Có thể dùng epsilon trong
hệ thống nâng cao, nhưng bản này tuân theo rule zero-variance của dự án.

\subsection{33. Code dùng population hay sample standard
deviation?}\label{33-code-duxf9ng-population-hay-sample-standard-deviation}

Population standard deviation, variance chia N. Rolling window được xem
là toàn bộ baseline runtime hiện tại, không phải mẫu nghiên cứu cần ước
lượng tổng thể bằng N-1.

\subsection{34. Buffer có trộn dữ liệu giữa hai thiết bị
không?}\label{34-buffer-cuxf3-trux1ed9n-dux1eef-liux1ec7u-giux1eefa-hai-thiux1ebft-bux1ecb-khuxf4ng}

Không. Key là \texttt{deviceId:sensorType}, ví dụ
\texttt{room-a:temperature} khác \texttt{room-b:temperature} và khác
\texttt{room-a:humidity}.

\subsection{35. Automatic mode có tắt hard threshold
không?}\label{35-automatic-mode-cuxf3-tux1eaft-hard-threshold-khuxf4ng}

Không. Hard threshold luôn được ưu tiên như safety bound. Automatic mode
chỉ bật thêm nhánh Z-score khi giá trị chưa vi phạm hard bound.

\subsection{36. Nhiệt độ Z-score âm có bật fan
không?}\label{36-nhiux1ec7t-ux111ux1ed9-z-score-uxe2m-cuxf3-bux1eadt-fan-khuxf4ng}

Không. Temperature anomaly vẫn được log và email nếu HIGH, nhưng
\texttt{autoFan} chỉ true khi giá trị mới lớn hơn mean, tức bất thường
theo hướng nóng lên.

\subsection{37. Gas vượt ngưỡng tạo severity
gì?}\label{37-gas-vux1b0ux1ee3t-ngux1b0ux1ee1ng-tux1ea1o-severity-guxec}

CRITICAL, trigger type \texttt{HARD\_THRESHOLD\_BREACH}. Edge alert từ
ESP32 dùng trigger type \texttt{EDGE\_SAFETY\_TRIGGER} nhưng cũng
CRITICAL.

\subsection{38. Motion tạo anomaly khi
nào?}\label{38-motion-tux1ea1o-anomaly-khi-nuxe0o}

Chỉ khi \texttt{motionDetected=true} và system đang armed. Nó tạo HIGH
anomaly, gửi email và publish lệnh bật alert LED.

\subsection{39. Humidity anomaly gửi email
không?}\label{39-humidity-anomaly-gux1eedi-email-khuxf4ng}

Không trong matrix hiện tại vì severity MEDIUM. Nó vẫn được lưu vào
\texttt{anomaly\_logs} và hiện trên dashboard.

\subsection{40. Dedupe key gồm gì?}\label{40-dedupe-key-gux1ed3m-guxec}

\texttt{deviceId:sensorType:triggerType}. Như vậy hard threshold và edge
trigger của cùng gas vẫn là hai loại sự kiện khác nhau.

\subsection{41. Vì sao reserve dedupe key trước action
async?}\label{41-vuxec-sao-reserve-dedupe-key-trux1b0ux1edbc-action-async}

Nếu hai telemetry message vào gần như đồng thời, cả hai có thể cùng vượt
kiểm tra trước khi lần đầu ghi DB. Reserve trước giúp chặn race
duplicate.

\subsection{42. Nếu action hoặc ghi DB lỗi thì dedupe key thế
nào?}\label{42-nux1ebfu-action-houx1eb7c-ghi-db-lux1ed7i-thuxec-dedupe-key-thux1ebf-nuxe0o}

Catch sẽ xóa reservation nếu nó vẫn thuộc lần xử lý hiện tại, cho phép
message sau retry thay vì bị chặn oan.

\subsection{43. Các anomaly có luôn được tạo đúng thứ tự gas, temp,
motion, humidity
không?}\label{43-cuxe1c-anomaly-cuxf3-luuxf4n-ux111ux1b0ux1ee3c-tux1ea1o-ux111uxfang-thux1ee9-tux1ef1-gas-temp-motion-humidity-khuxf4ng}

Không nên phụ thuộc thứ tự hoàn tất vì dùng \texttt{Promise.all}.
Business correctness nằm ở tập kết quả và nội dung từng anomaly.

\subsection{44. Nếu bật fan thất bại thì anomaly có bị mất
không?}\label{44-nux1ebfu-bux1eadt-fan-thux1ea5t-bux1ea1i-thuxec-anomaly-cuxf3-bux1ecb-mux1ea5t-khuxf4ng}

Không. Lỗi action được bắt và ghi vào \texttt{actionsTaken}; pipeline
tiếp tục tạo anomaly log. Đây là graceful degradation.

\subsection{45. Điểm hạn chế của rolling buffer hiện tại là
gì?}\label{45-ux111iux1ec3m-hux1ea1n-chux1ebf-cux1ee7a-rolling-buffer-hiux1ec7n-tux1ea1i-luxe0-guxec}

Buffer nằm trong RAM nên mất khi backend restart và mỗi instance có
buffer riêng. Production nhiều instance nên lưu state vào Redis hoặc
dùng truy vấn window từ time-series database.

\section{D. NC-03 - MongoDB Time
Series}\label{d-nc-03---mongodb-time-series}

\subsection{46. Collection telemetry tên
gì?}\label{46-collection-telemetry-tuxean-guxec}

\texttt{sensor\_readings}.

\subsection{\texorpdfstring{47. \texttt{timeField} và \texttt{metaField}
là
gì?}{47. timeField và metaField là gì?}}\label{47-timefield-vuxe0-metafield-luxe0-guxec}

\texttt{timestamp} là timeField; \texttt{deviceId} là metaField. MongoDB
dùng timestamp để tổ chức theo thời gian và deviceId để nhóm metadata
series.

\subsection{48. Granularity là gì?}\label{48-granularity-luxe0-guxec}

\texttt{seconds}, phù hợp dữ liệu ESP32 gửi mỗi khoảng 2 giây.

\subsection{49. Retention là bao
lâu?}\label{49-retention-luxe0-bao-luxe2u}

2,592,000 giây, tương đương 30 ngày.

\subsection{50. TTL có xóa đúng ngay khi đủ 30 ngày
không?}\label{50-ttl-cuxf3-xuxf3a-ux111uxfang-ngay-khi-ux111ux1ee7-30-nguxe0y-khuxf4ng}

Không chính xác từng giây vì TTL monitor chạy nền. Document sẽ được xóa
tự động sau khi hết hạn với độ trễ vận hành nhỏ.

\subsection{\texorpdfstring{51. Vì sao Mongoose schema dùng
\texttt{autoCreate=false}?}{51. Vì sao Mongoose schema dùng autoCreate=false?}}\label{51-vuxec-sao-mongoose-schema-duxf9ng-autocreatefalse}

Để Mongoose không tạo nhầm collection thường. Backend chủ động tạo Time
Series collection với đúng options trước khi ghi.

\subsection{52. Startup kiểm tra collection ra
sao?}\label{52-startup-kiux1ec3m-tra-collection-ra-sao}

\texttt{ensureSensorReadingsTimeSeries()} đọc metadata. Nếu chưa có thì
tạo; nếu có thì so timeField, metaField, granularity và TTL. Sai bất kỳ
option nào thì fail fast.

\subsection{53. Vì sao fail fast thay vì tự sửa
collection?}\label{53-vuxec-sao-fail-fast-thay-vuxec-tux1ef1-sux1eeda-collection}

Đổi một collection thường thành Time Series không phải update option đơn
giản và có nguy cơ mất dữ liệu. Fail fast buộc migration có chủ đích.

\subsection{54. Telemetry được validate trước hay sau khi
lưu?}\label{54-telemetry-ux111ux1b0ux1ee3c-validate-trux1b0ux1edbc-hay-sau-khi-lux1b0u}

Trước. \texttt{validateTelemetryPayload()} chuẩn hóa dữ liệu, kiểm tra
type/range và đối chiếu deviceId trong payload với topic.

\subsection{55. Vì sao phải đối chiếu deviceId payload với
topic?}\label{55-vuxec-sao-phux1ea3i-ux111ux1ed1i-chiux1ebfu-deviceid-payload-vux1edbi-topic}

Để client không publish vào topic của device A nhưng khai payload device
B, gây ghi sai series hoặc giả mạo dữ liệu.

\subsection{56. Tại sao historical query sort DESC rồi
reverse?}\label{56-tux1ea1i-sao-historical-query-sort-desc-rux1ed3i-reverse}

DESC + limit lấy đúng N mẫu mới nhất. Reverse sau đó trả chúng theo thời
gian tăng dần để chart vẽ đúng. ASC + limit sẽ lấy N mẫu cũ nhất.

\subsection{57. Historical API giới hạn bao nhiêu
record?}\label{57-historical-api-giux1edbi-hux1ea1n-bao-nhiuxeau-record}

\texttt{limit} phải là integer từ 1 đến 1000. Giới hạn ngăn response quá
lớn và query thiếu kiểm soát.

\subsection{58. Nếu chỉ hỏi một sensor thì API làm
gì?}\label{58-nux1ebfu-chux1ec9-hux1ecfi-mux1ed9t-sensor-thuxec-api-luxe0m-guxec}

Route dùng allowlist map sensor sang path schema rồi \texttt{.select()}
timestamp, deviceId và field sensor đó. Không lấy field tùy ý từ query
string.

\subsection{59. Dữ liệu được lưu trước hay phân tích anomaly
trước?}\label{59-dux1eef-liux1ec7u-ux111ux1b0ux1ee3c-lux1b0u-trux1b0ux1edbc-hay-phuxe2n-tuxedch-anomaly-trux1b0ux1edbc}

Lưu trước, sau đó chạy anomaly engine trên reading đã tạo. Như vậy dữ
liệu gốc vẫn có trong DB nếu action phía sau gặp lỗi.

\section{E. NC-05 - Email}\label{e-nc-05---email}

\subsection{60. Thư viện gửi email là
gì?}\label{60-thux1b0-viux1ec7n-gux1eedi-email-luxe0-guxec}

Nodemailer với SMTP transport.

\subsection{61. Severity nào gửi
email?}\label{61-severity-nuxe0o-gux1eedi-email}

HIGH và CRITICAL. LOW/MEDIUM trả status \texttt{ignored} trước khi liên
hệ SMTP.

\subsection{62. Throttle window là bao
lâu?}\label{62-throttle-window-luxe0-bao-luxe2u}

15 phút, tức \texttt{15\ *\ 60\ *\ 1000} millisecond.

\subsection{63. Throttle key gồm
gì?}\label{63-throttle-key-gux1ed3m-guxec}

\texttt{deviceId:sensorType}. Nó bảo đảm các cảm biến và thiết bị khác
nhau không chặn email của nhau.

\subsection{64. Vì sao key không chứa
triggerType?}\label{64-vuxec-sao-key-khuxf4ng-chux1ee9a-triggertype}

Mục tiêu là hạn chế email theo vấn đề cảm biến. Nếu cùng gas vừa có hard
threshold vừa có edge alert, người dùng vẫn không nên nhận hai chuỗi
email riêng dồn dập.

\subsection{65. Thời điểm throttle được ghi khi
nào?}\label{65-thux1eddi-ux111iux1ec3m-throttle-ux111ux1b0ux1ee3c-ghi-khi-nuxe0o}

Sau khi \texttt{sendMail()} thành công. SMTP lỗi không được phép mở một
window giả.

\subsection{66. Nếu email bị disable thì anomaly có mất
không?}\label{66-nux1ebfu-email-bux1ecb-disable-thuxec-anomaly-cuxf3-mux1ea5t-khuxf4ng}

Không. Email service trả status \texttt{disabled}; anomaly engine ghi
\texttt{EMAIL\_DISABLED} trong actionsTaken rồi vẫn lưu anomaly.

\subsection{67. SMTP password nằm ở
đâu?}\label{67-smtp-password-nux1eb1m-ux1edf-ux111uxe2u}

Trong environment variable, không commit vào Git. \texttt{.env.example}
chỉ chứa placeholder.

Địa chỉ gửi (\texttt{SMTP\_USER}) vẫn do server quản lý, nhưng địa chỉ
nhận lấy từ \texttt{User.email} đã đăng ký. Người dùng bật hoặc tắt nhận
mail tại Settings.

\subsection{68. Email có chống HTML injection
không?}\label{68-email-cuxf3-chux1ed1ng-html-injection-khuxf4ng}

Có. Các field động được escape \texttt{\&}, \texttt{\textless{}},
\texttt{\textgreater{}}, dấu nháy trước khi chèn vào HTML template.

\subsection{69. Transporter được tạo mỗi email
không?}\label{69-transporter-ux111ux1b0ux1ee3c-tux1ea1o-mux1ed7i-email-khuxf4ng}

Không. Service cache transporter trong runtime để tái sử dụng kết
nối/cấu hình.

\subsection{70. Email recipient hiện hỗ trợ nhiều user
không?}\label{70-email-recipient-hiux1ec7n-hux1ed7-trux1ee3-nhiux1ec1u-user-khuxf4ng}

Có ở mức tài khoản: backend lấy email profile của mọi user có
\texttt{settings.emailEnabled=true}, loại trùng địa chỉ và dùng BCC khi
có nhiều người nhận. Tuy nhiên hệ thống chưa ánh xạ device với owner,
nên mọi user đang bật đều nhận anomaly của mọi thiết bị.
\texttt{ALERT\_RECIPIENT\_EMAIL} chỉ fallback khi database chưa có user.

\subsection{71. Throttle có tồn tại sau backend restart
không?}\label{71-throttle-cuxf3-tux1ed3n-tux1ea1i-sau-backend-restart-khuxf4ng}

Không, vì map nằm trong RAM. Muốn bền vững và dùng nhiều instance cần
Redis hoặc collection riêng với atomic expiration.

\subsection{72. Tại sao test dùng fake transport và fake
clock?}\label{72-tux1ea1i-sao-test-duxf9ng-fake-transport-vuxe0-fake-clock}

Để test nhanh, lặp lại được và không gửi email thật. Fake clock cho phép
nhảy chính xác 15 phút mà không phải chờ thực tế.

\section{F. Câu hỏi phản biện và cách trả lời trung
thực}\label{f-cuxe2u-hux1ecfi-phux1ea3n-biux1ec7n-vuxe0-cuxe1ch-trux1ea3-lux1eddi-trung-thux1ef1c}

\subsection{73. Dự án đã hoàn toàn production-ready
chưa?}\label{73-dux1ef1-uxe1n-ux111uxe3-houxe0n-touxe0n-production-ready-chux1b0a}

Chưa. Core coursework và automated logic đã hoàn thành. Production cần
persistent distributed state, per-user authorization/device ownership,
observability, retry policy và end-to-end load/security testing.

\subsection{74. Test pass có chứng minh relay dưới 500 ms
không?}\label{74-test-pass-cuxf3-chux1ee9ng-minh-relay-dux1b0ux1edbi-500-ms-khuxf4ng}

Không. Simulation chỉ chứng minh contract phần mềm. TC-06 phải đo từ
thao tác web đến relay thật bằng log timestamp hoặc thiết bị đo.

\subsection{75. Test email có chứng minh nhận mail dưới 5 giây
không?}\label{75-test-email-cuxf3-chux1ee9ng-minh-nhux1eadn-mail-dux1b0ux1edbi-5-giuxe2y-khuxf4ng}

Không. Fake SMTP chứng minh logic severity và throttle. TC-05 với SMTP
thật mới đo latency giao mail.

\subsection{76. Tại sao không dùng machine learning phức tạp
hơn?}\label{76-tux1ea1i-sao-khuxf4ng-duxf9ng-machine-learning-phux1ee9c-tux1ea1p-hux1a1n}

Dataset nhỏ, thời gian dự án ngắn và yêu cầu cần giải thích được.
Rolling Z-score minh bạch, nhẹ, dễ test và phù hợp ESP32/IoT coursework
hơn một model khó hiệu chỉnh.

\subsection{77. Tại sao không chạy Z-score trên
ESP32?}\label{77-tux1ea1i-sao-khuxf4ng-chux1ea1y-z-score-truxean-esp32}

Backend đã nhận và lưu lịch sử, có tài nguyên tốt hơn và dễ cập nhật
threshold. ESP32 chỉ giữ logic an toàn bắt buộc phải hoạt động offline
như gas threshold.

\subsection{78. Nếu MongoDB lỗi thì email/anomaly có chạy
không?}\label{78-nux1ebfu-mongodb-lux1ed7i-thuxec-emailanomaly-cuxf3-chux1ea1y-khuxf4ng}

Trong pipeline hiện tại \texttt{SensorReading.create()} chạy trước
anomaly engine, nên ghi telemetry lỗi sẽ dừng xử lý anomaly của message
đó. Đây là tradeoff bảo đảm anomaly dựa trên dữ liệu đã persist. Edge
gas vẫn an toàn cục bộ trên ESP32.

\subsection{79. Nếu email đã gửi nhưng ghi anomaly log lỗi thì
sao?}\label{79-nux1ebfu-email-ux111uxe3-gux1eedi-nhux1b0ng-ghi-anomaly-log-lux1ed7i-thuxec-sao}

Action được thực hiện trước khi \texttt{AnomalyLog.create()} để lưu đầy
đủ \texttt{actionsTaken}. Vì vậy có khả năng email đã gửi nhưng DB write
sau đó lỗi. Thiết kế production tốt hơn có thể tạo log trạng thái
PENDING trước rồi update action hoặc dùng outbox pattern.

\subsection{80. Em sẽ mở file nào đầu tiên nếu được yêu cầu show
code?}\label{80-em-sux1ebd-mux1edf-file-nuxe0o-ux111ux1ea7u-tiuxean-nux1ebfu-ux111ux1b0ux1ee3c-yuxeau-cux1ea7u-show-code}

\begin{itemize}
\tightlist
\item
  CB-03: \nolinkurl{backend/routes/control.js}, sau đó
  \texttt{mqttService.js},
  \nolinkurl{firmware/src/mqtt/mqtt_client.cpp}, \texttt{main.cpp} và
  \nolinkurl{actuators/relay.cpp}.
\item
  NC-01: \nolinkurl{backend/services/anomalyEngine.js} tại
  \texttt{evaluateZScore} và \texttt{processTelemetryForAnomalies}.
\item
  NC-03: \texttt{timeSeriesService.js}, \texttt{telemetryService.js},
  rồi \nolinkurl{routes/telemetry.js}.
\item
  NC-05: \texttt{emailService.js} tại \texttt{sendAlertEmail}.
\end{itemize}

\section{Câu hỏi bổ sung về MQ2 sau hiệu
chỉnh}\label{cuxe2u-hux1ecfi-bux1ed5-sung-vux1ec1-mq2-sau-hiux1ec7u-chux1ec9nh}

\subsection{\texorpdfstring{81. Tại sao \texttt{gasLevel} có thể là
null?}{81. Tại sao gasLevel có thể là null?}}\label{81-tux1ea1i-sao-gaslevel-cuxf3-thux1ec3-luxe0-null}

Vì MQ2 chưa ở trạng thái \texttt{ready}. \texttt{null} biểu diễn đúng
việc PPM chưa có giá trị, khác với \texttt{0} vốn là một phép đo.
Backend vẫn lưu \texttt{sensorStatus.gas} nhưng không tạo gas anomaly
hoặc email từ dữ liệu chưa hiệu chuẩn.

\subsection{\texorpdfstring{82. Vì sao cần burn-in 48 giờ và lưu
\texttt{Ro} trong
NVS?}{82. Vì sao cần burn-in 48 giờ và lưu Ro trong NVS?}}\label{82-vuxec-sao-cux1ea7n-burn-in-48-giux1edd-vuxe0-lux1b0u-ro-trong-nvs}

Datasheet MQ2 yêu cầu preheat ban đầu trên 48 giờ. \texttt{Ro} là điện
trở cảm biến trong không khí sạch; hiệu chuẩn ngay khi bật nguồn sẽ học
một baseline sai. NVS giữ \texttt{Ro} qua reset, nên các lần sau chỉ
warm-up 5 phút và không hiệu chuẩn sai lại từ đầu.

\subsection{83. Trong 48 giờ burn-in có mất an toàn local
không?}\label{83-trong-48-giux1edd-burn-in-cuxf3-mux1ea5t-an-touxe0n-local-khuxf4ng}

Không hoàn toàn. Firmware chưa được phép gọi điện áp raw là PPM hoặc gửi
email, nhưng vẫn có ngưỡng ADC gần bão hòa để bật còi, LED và OLED tại
chỗ. Đây là lớp dự phòng bảo thủ, không thay cho cảnh báo 600 PPM sau
hiệu chuẩn.

\subsection{84. Tại sao phải dùng FreeRTOS task riêng cho
gas?}\label{84-tux1ea1i-sao-phux1ea3i-duxf9ng-freertos-task-riuxeang-cho-gas}

WiFi và bắt tay TLS có thể chặn Arduino \texttt{loop()} lâu hơn 200 ms.
Task gas ưu tiên cao chạy mỗi 100 ms, chỉ đọc ADC và điều khiển local;
task mạng publish bản ghi đã xếp hàng sau. Vì vậy đường an toàn không
chờ cloud.

\subsection{85. MQ2 có đo PPM chính xác cho mọi loại khí
không?}\label{85-mq2-cuxf3-ux111o-ppm-chuxednh-xuxe1c-cho-mux1ecdi-loux1ea1i-khuxed-khuxf4ng}

Không. MQ2 phản ứng chéo với propane, methane, cồn và khói. Hệ thống nội
suy đường propane trong datasheet nên phải gọi là LPG-equivalent PPM
estimate. Muốn PPM chuẩn cần khí chuẩn, điều kiện nhiệt độ/độ ẩm kiểm
soát và cảm biến chọn lọc phù hợp.

\chapter{Code walkthrough phần Thái
Kiệt}\label{code-walkthrough-phux1ea7n-thuxe1i-kiux1ec7t}

Tài liệu này không sao chép toàn bộ source. Nó trích các đoạn quyết định
hành vi của CB-03, NC-01, NC-03 và NC-05, rồi giải thích cách các đoạn
nối với nhau. Khi vấn đáp, nên mở file thật và tìm theo tên hàm in đậm.

\section{1. Cách đọc project}\label{1-cuxe1ch-ux111ux1ecdc-project}

Đọc một luồng theo thứ tự sau:

\begin{Shaded}
\begin{Highlighting}[]
\NormalTok{Input {-}\textgreater{} route/message handler {-}\textgreater{} service {-}\textgreater{} model/output {-}\textgreater{} test}
\end{Highlighting}
\end{Shaded}

Ví dụ CB-03:

\begin{Shaded}
\begin{Highlighting}[]
\NormalTok{Overview.jsx}
\NormalTok{  {-}\textgreater{} POST /api/v1/control/actuator}
\NormalTok{  {-}\textgreater{} control.js}
\NormalTok{  {-}\textgreater{} mqttService.publishControlCommand()}
\NormalTok{  {-}\textgreater{} MQTT topic}
\NormalTok{  {-}\textgreater{} MqttClient.handleMessage()}
\NormalTok{  {-}\textgreater{} main.cpp onControlMessage()}
\NormalTok{  {-}\textgreater{} Relay.on()/off()}
\NormalTok{  {-}\textgreater{} digitalWrite(GPIO 16 theo RELAY\_ACTIVE\_LOW)}
\end{Highlighting}
\end{Shaded}

\section{2. CB-03 - Điều khiển
quạt}\label{2-cb-03---ux111iux1ec1u-khiux1ec3n-quux1ea1t}

\subsection{2.1 Frontend gửi lệnh}\label{21-frontend-gux1eedi-lux1ec7nh}

File:
\href{../frontend/src/pages/Overview.jsx}{\nolinkurl{frontend/src/pages/Overview.jsx}}

Đầu file có mapping giữa tên actuator của API và field trong telemetry:

\begin{Shaded}
\begin{Highlighting}[]
\KeywordTok{const}\NormalTok{ ACTUATOR\_STATE\_FIELDS }\OperatorTok{=}\NormalTok{ \{}
  \DataTypeTok{fan}\OperatorTok{:} \StringTok{\textquotesingle{}fanOn\textquotesingle{}}\OperatorTok{,}
  \DataTypeTok{buzzer}\OperatorTok{:} \StringTok{\textquotesingle{}buzzerActive\textquotesingle{}}\OperatorTok{,}
  \DataTypeTok{alertLed}\OperatorTok{:} \StringTok{\textquotesingle{}alertLedActive\textquotesingle{}}
\NormalTok{\}}\OperatorTok{;}
\end{Highlighting}
\end{Shaded}

Mapping này thay cho ternary lồng nhau. Khi thêm actuator mới, chỉ cần
thêm một entry và phần cập nhật optimistic state vẫn dễ đọc.

Hàm \textbf{handleActuatorToggle} thực hiện ba việc:

\begin{Shaded}
\begin{Highlighting}[]
\KeywordTok{const}\NormalTok{ targetState }\OperatorTok{=} \OperatorTok{!}\NormalTok{currentState}\OperatorTok{;}
\KeywordTok{const}\NormalTok{ res }\OperatorTok{=} \ControlFlowTok{await} \FunctionTok{apiFetch}\NormalTok{(}\StringTok{\textquotesingle{}/api/v1/control/actuator\textquotesingle{}}\OperatorTok{,}\NormalTok{ \{}
  \DataTypeTok{method}\OperatorTok{:} \StringTok{\textquotesingle{}POST\textquotesingle{}}\OperatorTok{,}
  \DataTypeTok{headers}\OperatorTok{:}\NormalTok{ \{ }\StringTok{\textquotesingle{}Content{-}Type\textquotesingle{}}\OperatorTok{:} \StringTok{\textquotesingle{}application/json\textquotesingle{}}\NormalTok{ \}}\OperatorTok{,}
  \DataTypeTok{body}\OperatorTok{:} \BuiltInTok{JSON}\OperatorTok{.}\FunctionTok{stringify}\NormalTok{(\{}
    \DataTypeTok{deviceId}\OperatorTok{:}\NormalTok{ telemetry}\OperatorTok{?.}\AttributeTok{deviceId} \OperatorTok{||} \StringTok{\textquotesingle{}safebox{-}01\textquotesingle{}}\OperatorTok{,}
\NormalTok{    actuator}\OperatorTok{,}
    \DataTypeTok{state}\OperatorTok{:}\NormalTok{ targetState}
\NormalTok{  \})}
\NormalTok{\})}\OperatorTok{;}
\end{Highlighting}
\end{Shaded}

\begin{enumerate}
\def\labelenumi{\arabic{enumi}.}
\tightlist
\item
  Đảo trạng thái hiện tại.
\item
  Gửi JSON đến backend.
\item
  Dùng \texttt{apiFetch} để tự gắn JWT.
\end{enumerate}

Sau response thành công, UI cập nhật trước khi vòng polling tiếp theo
chạy:

\begin{Shaded}
\begin{Highlighting}[]
\KeywordTok{const}\NormalTok{ stateField }\OperatorTok{=}\NormalTok{ ACTUATOR\_STATE\_FIELDS[actuator]}\OperatorTok{;}
\ControlFlowTok{if}\NormalTok{ (}\OperatorTok{!}\NormalTok{stateField) }\ControlFlowTok{throw} \KeywordTok{new} \BuiltInTok{Error}\NormalTok{(}\VerbatimStringTok{\textasciigrave{}Unsupported actuator \textquotesingle{}}\SpecialCharTok{$\{}\NormalTok{actuator}\SpecialCharTok{\}}\VerbatimStringTok{\textquotesingle{}\textasciigrave{}}\NormalTok{)}\OperatorTok{;}
\FunctionTok{setTelemetry}\NormalTok{((prev) }\KeywordTok{=\textgreater{}}\NormalTok{ (\{}
  \OperatorTok{...}\NormalTok{(prev }\OperatorTok{||}\NormalTok{ \{\})}\OperatorTok{,}
  \DataTypeTok{actuatorStates}\OperatorTok{:}\NormalTok{ \{}
    \OperatorTok{...}\NormalTok{(prev}\OperatorTok{?.}\AttributeTok{actuatorStates} \OperatorTok{||}\NormalTok{ \{\})}\OperatorTok{,}
\NormalTok{    [stateField]}\OperatorTok{:}\NormalTok{ targetState}
\NormalTok{  \}}
\NormalTok{\}))}\OperatorTok{;}
\end{Highlighting}
\end{Shaded}

Đây là optimistic update. Nó giúp UI phản hồi nhanh nhưng không được xem
là xác nhận phần cứng.

\subsection{2.2 JWT helper}\label{22-jwt-helper}

File: \href{../frontend/src/api.js}{\nolinkurl{frontend/src/api.js}}

\begin{Shaded}
\begin{Highlighting}[]
\ImportTok{export} \KeywordTok{function} \FunctionTok{apiFetch}\NormalTok{(input}\OperatorTok{,}\NormalTok{ options }\OperatorTok{=}\NormalTok{ \{\}) \{}
  \KeywordTok{const}\NormalTok{ token }\OperatorTok{=}\NormalTok{ localStorage}\OperatorTok{.}\FunctionTok{getItem}\NormalTok{(}\StringTok{\textquotesingle{}safebox\_token\textquotesingle{}}\NormalTok{)}\OperatorTok{;}
  \KeywordTok{const}\NormalTok{ headers }\OperatorTok{=} \KeywordTok{new} \FunctionTok{Headers}\NormalTok{(options}\OperatorTok{.}\AttributeTok{headers} \OperatorTok{||}\NormalTok{ \{\})}\OperatorTok{;}

  \ControlFlowTok{if}\NormalTok{ (token) headers}\OperatorTok{.}\FunctionTok{set}\NormalTok{(}\StringTok{\textquotesingle{}Authorization\textquotesingle{}}\OperatorTok{,} \VerbatimStringTok{\textasciigrave{}Bearer }\SpecialCharTok{$\{}\NormalTok{token}\SpecialCharTok{\}}\VerbatimStringTok{\textasciigrave{}}\NormalTok{)}\OperatorTok{;}

  \ControlFlowTok{return} \FunctionTok{fetch}\NormalTok{(input}\OperatorTok{,}\NormalTok{ \{}
    \OperatorTok{...}\NormalTok{options}\OperatorTok{,}
\NormalTok{    headers}
\NormalTok{  \})}\OperatorTok{;}
\NormalTok{\}}
\end{Highlighting}
\end{Shaded}

Điểm cần nói:

\begin{itemize}
\tightlist
\item
  Helper không tự tạo token; token được tạo bởi backend khi login.
\item
  Nó giữ nguyên header caller truyền vào rồi bổ sung Authorization.
\item
  Nếu không có token, backend middleware trả 401.
\end{itemize}

\subsection{2.3 REST route nhận
lệnh}\label{23-rest-route-nhux1eadn-lux1ec7nh}

File:
\href{../backend/routes/control.js}{\nolinkurl{backend/routes/control.js}}

Hàm \textbf{parseActuatorState} chuẩn hóa hai dạng input:

\begin{Shaded}
\begin{Highlighting}[]
\KeywordTok{function} \FunctionTok{parseActuatorState}\NormalTok{(value) \{}
  \ControlFlowTok{if}\NormalTok{ (}\KeywordTok{typeof}\NormalTok{ value }\OperatorTok{===} \StringTok{\textquotesingle{}boolean\textquotesingle{}}\NormalTok{) }\ControlFlowTok{return}\NormalTok{ value}\OperatorTok{;}
  \ControlFlowTok{if}\NormalTok{ (value }\OperatorTok{===} \StringTok{\textquotesingle{}on\textquotesingle{}}\NormalTok{) }\ControlFlowTok{return} \KeywordTok{true}\OperatorTok{;}
  \ControlFlowTok{if}\NormalTok{ (value }\OperatorTok{===} \StringTok{\textquotesingle{}off\textquotesingle{}}\NormalTok{) }\ControlFlowTok{return} \KeywordTok{false}\OperatorTok{;}
  \ControlFlowTok{throw} \KeywordTok{new} \BuiltInTok{TypeError}\NormalTok{(}\StringTok{"state must be a boolean or the string \textquotesingle{}on\textquotesingle{}/\textquotesingle{}off\textquotesingle{}"}\NormalTok{)}\OperatorTok{;}
\NormalTok{\}}
\end{Highlighting}
\end{Shaded}

Route kiểm tra actuator trước khi publish:

\begin{Shaded}
\begin{Highlighting}[]
\ControlFlowTok{if}\NormalTok{ (}\OperatorTok{!}\NormalTok{ALLOWED\_ACTUATORS}\OperatorTok{.}\FunctionTok{has}\NormalTok{(actuator)) \{}
  \ControlFlowTok{return}\NormalTok{ res}\OperatorTok{.}\FunctionTok{status}\NormalTok{(}\DecValTok{400}\NormalTok{)}\OperatorTok{.}\FunctionTok{json}\NormalTok{(\{}
    \DataTypeTok{success}\OperatorTok{:} \KeywordTok{false}\OperatorTok{,}
    \DataTypeTok{error}\OperatorTok{:} \VerbatimStringTok{\textasciigrave{}Unsupported actuator \textquotesingle{}}\SpecialCharTok{$\{}\NormalTok{actuator}\SpecialCharTok{\}}\VerbatimStringTok{\textquotesingle{}\textasciigrave{}}
\NormalTok{  \})}\OperatorTok{;}
\NormalTok{\}}

\KeywordTok{const}\NormalTok{ targetState }\OperatorTok{=} \FunctionTok{parseActuatorState}\NormalTok{(state)}\OperatorTok{;}
\KeywordTok{const}\NormalTok{ publishResult }\OperatorTok{=} \ControlFlowTok{await} \FunctionTok{publishControlCommand}\NormalTok{(deviceId}\OperatorTok{,}\NormalTok{ actuator}\OperatorTok{,}\NormalTok{ targetState)}\OperatorTok{;}
\end{Highlighting}
\end{Shaded}

Nếu MQTT thật đang offline, service gắn code \texttt{MQTT\_OFFLINE} và
route trả 503. Input sai trả 400. Lỗi không dự kiến trả 500.

\subsection{2.4 Kho trạng thái
actuator}\label{24-kho-trux1ea1ng-thuxe1i-actuator}

File:
\href{../backend/services/actuatorStateService.js}{\nolinkurl{backend/services/actuatorStateService.js}}

\begin{Shaded}
\begin{Highlighting}[]
\KeywordTok{const}\NormalTok{ ACTUATOR\_STATE\_FIELDS }\OperatorTok{=} \BuiltInTok{Object}\OperatorTok{.}\FunctionTok{freeze}\NormalTok{(\{}
  \DataTypeTok{fan}\OperatorTok{:} \StringTok{\textquotesingle{}fanOn\textquotesingle{}}\OperatorTok{,}
  \DataTypeTok{buzzer}\OperatorTok{:} \StringTok{\textquotesingle{}buzzerActive\textquotesingle{}}\OperatorTok{,}
  \DataTypeTok{alertLed}\OperatorTok{:} \StringTok{\textquotesingle{}alertLedActive\textquotesingle{}}
\NormalTok{\})}\OperatorTok{;}
\end{Highlighting}
\end{Shaded}

Hàm \textbf{updateActuatorState} lưu requested state sau publish:

\begin{Shaded}
\begin{Highlighting}[]
\KeywordTok{function} \FunctionTok{updateActuatorState}\NormalTok{(actuator}\OperatorTok{,}\NormalTok{ state) \{}
  \KeywordTok{const}\NormalTok{ stateField }\OperatorTok{=}\NormalTok{ ACTUATOR\_STATE\_FIELDS[actuator]}\OperatorTok{;}
  \ControlFlowTok{if}\NormalTok{ (}\OperatorTok{!}\NormalTok{stateField) }\ControlFlowTok{throw} \KeywordTok{new} \BuiltInTok{TypeError}\NormalTok{(}\VerbatimStringTok{\textasciigrave{}Unsupported actuator \textquotesingle{}}\SpecialCharTok{$\{}\NormalTok{actuator}\SpecialCharTok{\}}\VerbatimStringTok{\textquotesingle{}\textasciigrave{}}\NormalTok{)}\OperatorTok{;}

\NormalTok{  currentActuatorStates }\OperatorTok{=}\NormalTok{ \{}
    \OperatorTok{...}\NormalTok{currentActuatorStates}\OperatorTok{,}
\NormalTok{    [stateField]}\OperatorTok{:} \BuiltInTok{Boolean}\NormalTok{(state)}
\NormalTok{  \}}\OperatorTok{;}
  \ControlFlowTok{return} \FunctionTok{getCurrentActuatorStates}\NormalTok{()}\OperatorTok{;}
\NormalTok{\}}
\end{Highlighting}
\end{Shaded}

Hàm \textbf{synchronizeActuatorStates} thay snapshot bằng telemetry đã
xác thực:

\begin{Shaded}
\begin{Highlighting}[]
\KeywordTok{function} \FunctionTok{synchronizeActuatorStates}\NormalTok{(actuatorStates) \{}
\NormalTok{  currentActuatorStates }\OperatorTok{=}\NormalTok{ \{}
    \DataTypeTok{fanOn}\OperatorTok{:} \BuiltInTok{Boolean}\NormalTok{(actuatorStates}\OperatorTok{.}\AttributeTok{fanOn}\NormalTok{)}\OperatorTok{,}
    \DataTypeTok{buzzerActive}\OperatorTok{:} \BuiltInTok{Boolean}\NormalTok{(actuatorStates}\OperatorTok{.}\AttributeTok{buzzerActive}\NormalTok{)}\OperatorTok{,}
    \DataTypeTok{alertLedActive}\OperatorTok{:} \BuiltInTok{Boolean}\NormalTok{(actuatorStates}\OperatorTok{.}\AttributeTok{alertLedActive}\NormalTok{)}
\NormalTok{  \}}\OperatorTok{;}
  \ControlFlowTok{return} \FunctionTok{getCurrentActuatorStates}\NormalTok{()}\OperatorTok{;}
\NormalTok{\}}
\end{Highlighting}
\end{Shaded}

Service trả object copy, không trả reference map nội bộ. Caller không
thể vô tình sửa state mà không đi qua service.

\subsection{2.5 Backend publish MQTT}\label{25-backend-publish-mqtt}

File:
\href{../backend/services/mqttService.js}{\nolinkurl{backend/services/mqttService.js}}

Hàm \textbf{publishControlCommand}:

\begin{Shaded}
\begin{Highlighting}[]
\KeywordTok{const}\NormalTok{ topic }\OperatorTok{=} \VerbatimStringTok{\textasciigrave{}safebox/}\SpecialCharTok{$\{}\NormalTok{deviceId}\SpecialCharTok{\}}\VerbatimStringTok{/control\textasciigrave{}}\OperatorTok{;}
\KeywordTok{const}\NormalTok{ payload }\OperatorTok{=} \BuiltInTok{JSON}\OperatorTok{.}\FunctionTok{stringify}\NormalTok{(\{ actuator}\OperatorTok{,} \DataTypeTok{state}\OperatorTok{:}\NormalTok{ state }\OperatorTok{?} \StringTok{\textquotesingle{}on\textquotesingle{}} \OperatorTok{:} \StringTok{\textquotesingle{}off\textquotesingle{}}\NormalTok{ \})}\OperatorTok{;}
\ControlFlowTok{await} \KeywordTok{new} \BuiltInTok{Promise}\NormalTok{((resolve}\OperatorTok{,}\NormalTok{ reject) }\KeywordTok{=\textgreater{}}\NormalTok{ \{}
\NormalTok{  mqttClient}\OperatorTok{.}\FunctionTok{publish}\NormalTok{(topic}\OperatorTok{,}\NormalTok{ payload}\OperatorTok{,}\NormalTok{ \{ }\DataTypeTok{qos}\OperatorTok{:} \DecValTok{1}\OperatorTok{,} \DataTypeTok{retain}\OperatorTok{:} \KeywordTok{false}\NormalTok{ \}}\OperatorTok{,}\NormalTok{ (error) }\KeywordTok{=\textgreater{}}\NormalTok{ \{}
    \ControlFlowTok{if}\NormalTok{ (error) }\FunctionTok{reject}\NormalTok{(error)}\OperatorTok{;}
    \ControlFlowTok{else} \FunctionTok{resolve}\NormalTok{()}\OperatorTok{;}
\NormalTok{  \})}\OperatorTok{;}
\NormalTok{\})}\OperatorTok{;}
\FunctionTok{updateActuatorState}\NormalTok{(actuator}\OperatorTok{,}\NormalTok{ state)}\OperatorTok{;}
\ControlFlowTok{return}\NormalTok{ \{ }\DataTypeTok{status}\OperatorTok{:} \StringTok{\textquotesingle{}published\textquotesingle{}}\OperatorTok{,}\NormalTok{ topic \}}\OperatorTok{;}
\end{Highlighting}
\end{Shaded}

Ý nghĩa từng option:

\begin{itemize}
\tightlist
\item
  \texttt{qos:\ 1}: at least once.
\item
  \texttt{retain:\ false}: lệnh cũ không được giữ lại để chạy lại khi
  thiết bị kết nối sau.
\item
  Chỉ cập nhật requested state sau callback publish không báo lỗi.
\end{itemize}

Khi telemetry thật quay về:

\begin{Shaded}
\begin{Highlighting}[]
\ControlFlowTok{if}\NormalTok{ (channel }\OperatorTok{===} \StringTok{\textquotesingle{}telemetry\textquotesingle{}}\NormalTok{) \{}
  \KeywordTok{const}\NormalTok{ result }\OperatorTok{=} \ControlFlowTok{await} \FunctionTok{ingestTelemetry}\NormalTok{(payload}\OperatorTok{,}\NormalTok{ deviceId)}\OperatorTok{;}
  \FunctionTok{synchronizeActuatorStates}\NormalTok{(result}\OperatorTok{.}\AttributeTok{reading}\OperatorTok{.}\AttributeTok{actuatorStates}\NormalTok{)}\OperatorTok{;}
  \ControlFlowTok{return}\NormalTok{ result}\OperatorTok{;}
\NormalTok{\}}
\end{Highlighting}
\end{Shaded}

\subsection{2.6 Firmware nhận MQTT}\label{26-firmware-nhux1eadn-mqtt}

File:
\href{../firmware/src/mqtt/mqtt_client.cpp}{\nolinkurl{firmware/src/mqtt/mqtt_client.cpp}}

\begin{Shaded}
\begin{Highlighting}[]
\DataTypeTok{void}\NormalTok{ MqttClient}\OperatorTok{::}\NormalTok{handleMessage}\OperatorTok{(}\NormalTok{byte}\OperatorTok{*}\NormalTok{ payload}\OperatorTok{,} \DataTypeTok{unsigned} \DataTypeTok{int}\NormalTok{ length}\OperatorTok{)} \OperatorTok{\{}
\NormalTok{  JsonDocument doc}\OperatorTok{;}
\NormalTok{  DeserializationError error }\OperatorTok{=}\NormalTok{ deserializeJson}\OperatorTok{(}\NormalTok{doc}\OperatorTok{,}\NormalTok{ payload}\OperatorTok{,}\NormalTok{ length}\OperatorTok{);}
  \ControlFlowTok{if} \OperatorTok{(}\NormalTok{error}\OperatorTok{)} \ControlFlowTok{return}\OperatorTok{;}

  \AttributeTok{const} \DataTypeTok{char}\OperatorTok{*}\NormalTok{ actuator }\OperatorTok{=}\NormalTok{ doc}\OperatorTok{[}\StringTok{"actuator"}\OperatorTok{]} \OperatorTok{|} \StringTok{""}\OperatorTok{;}
  \DataTypeTok{bool}\NormalTok{ state }\OperatorTok{=} \KeywordTok{false}\OperatorTok{;}
  \ControlFlowTok{if} \OperatorTok{(}\NormalTok{doc}\OperatorTok{[}\StringTok{"state"}\OperatorTok{].}\NormalTok{is}\OperatorTok{\textless{}}\DataTypeTok{bool}\OperatorTok{\textgreater{}())} \OperatorTok{\{}
\NormalTok{    state }\OperatorTok{=}\NormalTok{ doc}\OperatorTok{[}\StringTok{"state"}\OperatorTok{].}\NormalTok{as}\OperatorTok{\textless{}}\DataTypeTok{bool}\OperatorTok{\textgreater{}();}
  \OperatorTok{\}} \ControlFlowTok{else} \OperatorTok{\{}
    \AttributeTok{const} \DataTypeTok{char}\OperatorTok{*}\NormalTok{ stateStr }\OperatorTok{=}\NormalTok{ doc}\OperatorTok{[}\StringTok{"state"}\OperatorTok{]} \OperatorTok{|} \StringTok{"off"}\OperatorTok{;}
\NormalTok{    state }\OperatorTok{=} \OperatorTok{(}\NormalTok{strcmp}\OperatorTok{(}\NormalTok{stateStr}\OperatorTok{,} \StringTok{"on"}\OperatorTok{)} \OperatorTok{==} \DecValTok{0}\OperatorTok{);}
  \OperatorTok{\}}

  \ControlFlowTok{if} \OperatorTok{(}\NormalTok{\_controlCallback}\OperatorTok{)}\NormalTok{ \_controlCallback}\OperatorTok{(}\NormalTok{actuator}\OperatorTok{,}\NormalTok{ state}\OperatorTok{);}
\OperatorTok{\}}
\end{Highlighting}
\end{Shaded}

\texttt{MqttClient} chịu trách nhiệm subscribe topic control, parse JSON
và đổi state sang boolean. JSON lỗi bị bỏ qua. Chuỗi khác \texttt{on}
được xử lý theo hướng fail-off; actuator lạ không tạo tác động vì
callback ở \texttt{main.cpp} không có nhánh tương ứng.

File: \href{../firmware/src/main.cpp}{\nolinkurl{firmware/src/main.cpp}}

\begin{Shaded}
\begin{Highlighting}[]
\AttributeTok{static} \DataTypeTok{void}\NormalTok{ onControlMessage}\OperatorTok{(}\AttributeTok{const} \DataTypeTok{char}\OperatorTok{*}\NormalTok{ actuator}\OperatorTok{,} \DataTypeTok{bool}\NormalTok{ state}\OperatorTok{)} \OperatorTok{\{}
  \ControlFlowTok{if} \OperatorTok{(}\NormalTok{strcmp}\OperatorTok{(}\NormalTok{actuator}\OperatorTok{,} \StringTok{"fan"}\OperatorTok{)} \OperatorTok{==} \DecValTok{0}\OperatorTok{)} \OperatorTok{\{}
    \ControlFlowTok{if} \OperatorTok{(}\NormalTok{state}\OperatorTok{)}\NormalTok{ fanRelay}\OperatorTok{.}\NormalTok{on}\OperatorTok{();}
    \ControlFlowTok{else}\NormalTok{ fanRelay}\OperatorTok{.}\NormalTok{off}\OperatorTok{();}
\NormalTok{    DBG}\OperatorTok{(}\StringTok{"[Fan] Relay state: }\SpecialCharTok{\%s\textbackslash{}n}\StringTok{"}\OperatorTok{,}\NormalTok{ fanRelay}\OperatorTok{.}\NormalTok{isOn}\OperatorTok{()} \OperatorTok{?} \StringTok{"ON"} \OperatorTok{:} \StringTok{"OFF"}\OperatorTok{);}
  \OperatorTok{\}} \ControlFlowTok{else} \ControlFlowTok{if} \OperatorTok{(}\NormalTok{strcmp}\OperatorTok{(}\NormalTok{actuator}\OperatorTok{,} \StringTok{"buzzer"}\OperatorTok{)} \OperatorTok{==} \DecValTok{0}\OperatorTok{)} \OperatorTok{\{}
\NormalTok{    edgeSafety}\OperatorTok{.}\NormalTok{applyRemoteBuzzer}\OperatorTok{(}\NormalTok{state}\OperatorTok{);}
  \OperatorTok{\}} \ControlFlowTok{else} \ControlFlowTok{if} \OperatorTok{(}\NormalTok{strcmp}\OperatorTok{(}\NormalTok{actuator}\OperatorTok{,} \StringTok{"alertLed"}\OperatorTok{)} \OperatorTok{==} \DecValTok{0}\OperatorTok{)} \OperatorTok{\{}
\NormalTok{    edgeSafety}\OperatorTok{.}\NormalTok{applyRemoteAlertLed}\OperatorTok{(}\NormalTok{state}\OperatorTok{);}
  \OperatorTok{\}}
\OperatorTok{\}}
\end{Highlighting}
\end{Shaded}

\texttt{main.cpp} chỉ định tuyến actuator. Buzzer và alert LED đi qua
\texttt{EdgeSafety}, nhờ đó cảnh báo gas cục bộ có thể ưu tiên hơn lệnh
remote.

File:
\href{../firmware/src/actuators/relay.cpp}{\nolinkurl{firmware/src/actuators/relay.cpp}}

\begin{Shaded}
\begin{Highlighting}[]
\DataTypeTok{void}\NormalTok{ Relay}\OperatorTok{::}\NormalTok{begin}\OperatorTok{()} \OperatorTok{\{}
\NormalTok{  pinMode}\OperatorTok{(}\NormalTok{\_pin}\OperatorTok{,}\NormalTok{ OUTPUT}\OperatorTok{);}
\NormalTok{  off}\OperatorTok{();}
\OperatorTok{\}}

\DataTypeTok{void}\NormalTok{ Relay}\OperatorTok{::}\NormalTok{on}\OperatorTok{()} \OperatorTok{\{}
\NormalTok{  digitalWrite}\OperatorTok{(}\NormalTok{\_pin}\OperatorTok{,}\NormalTok{ RELAY\_ACTIVE\_LOW }\OperatorTok{?}\NormalTok{ LOW }\OperatorTok{:}\NormalTok{ HIGH}\OperatorTok{);}
\NormalTok{  \_state }\OperatorTok{=} \KeywordTok{true}\OperatorTok{;}
\OperatorTok{\}}

\DataTypeTok{void}\NormalTok{ Relay}\OperatorTok{::}\NormalTok{off}\OperatorTok{()} \OperatorTok{\{}
\NormalTok{  digitalWrite}\OperatorTok{(}\NormalTok{\_pin}\OperatorTok{,}\NormalTok{ RELAY\_ACTIVE\_LOW }\OperatorTok{?}\NormalTok{ HIGH }\OperatorTok{:}\NormalTok{ LOW}\OperatorTok{);}
\NormalTok{  \_state }\OperatorTok{=} \KeywordTok{false}\OperatorTok{;}
\OperatorTok{\}}
\end{Highlighting}
\end{Shaded}

Với module HL-52S mặc định \texttt{RELAY\_ACTIVE\_LOW=1}, GPIO 16 LOW là
relay ON và HIGH là OFF. \texttt{begin()} gọi \texttt{off()}, nên fan
vẫn tắt trong lúc boot.

\section{3. NC-01 - Z-score và rule
engine}\label{3-nc-01---z-score-vuxe0-rule-engine}

File:
\href{../backend/services/anomalyEngine.js}{\nolinkurl{backend/services/anomalyEngine.js}}

\subsection{3.1 Hằng số và
state}\label{31-hux1eb1ng-sux1ed1-vuxe0-state}

\begin{Shaded}
\begin{Highlighting}[]
\KeywordTok{const}\NormalTok{ WINDOW\_SIZE }\OperatorTok{=} \DecValTok{30}\OperatorTok{;}
\KeywordTok{const}\NormalTok{ DEDUPE\_WINDOW\_MS }\OperatorTok{=} \BuiltInTok{Number}\NormalTok{(}\BuiltInTok{process}\OperatorTok{.}\AttributeTok{env}\OperatorTok{.}\AttributeTok{ANOMALY\_DEDUPE\_WINDOW\_MS} \OperatorTok{||} \DecValTok{60\_000}\NormalTok{)}\OperatorTok{;}
\KeywordTok{const}\NormalTok{ MIN\_DELTA }\OperatorTok{=}\NormalTok{ \{ }\DataTypeTok{temperature}\OperatorTok{:} \FloatTok{1.0}\OperatorTok{,} \DataTypeTok{humidity}\OperatorTok{:} \FloatTok{2.0}\NormalTok{ \}}\OperatorTok{;}
\end{Highlighting}
\end{Shaded}

\texttt{DEFAULT\_THRESHOLDS} được freeze để tránh sửa nhầm object mặc
định. Runtime dùng một bản copy, nên \texttt{updateThresholds()} vẫn
thay đổi được cấu hình đang chạy.

\begin{Shaded}
\begin{Highlighting}[]
\KeywordTok{let}\NormalTok{ systemThresholds }\OperatorTok{=}\NormalTok{ \{ }\OperatorTok{...}\NormalTok{DEFAULT\_THRESHOLDS \}}\OperatorTok{;}
\KeywordTok{const}\NormalTok{ rollingBuffers }\OperatorTok{=} \KeywordTok{new} \BuiltInTok{Map}\NormalTok{()}\OperatorTok{;}
\KeywordTok{const}\NormalTok{ anomalyDedupMap }\OperatorTok{=} \KeywordTok{new} \BuiltInTok{Map}\NormalTok{()}\OperatorTok{;}
\end{Highlighting}
\end{Shaded}

Key rolling buffer là \texttt{deviceId:sensorType}, vì mỗi thiết bị và
mỗi metric cần baseline riêng.

\subsection{3.2 Tính mean và standard
deviation}\label{32-tuxednh-mean-vuxe0-standard-deviation}

\begin{Shaded}
\begin{Highlighting}[]
\KeywordTok{function} \FunctionTok{meanAndStdDev}\NormalTok{(values) \{}
  \KeywordTok{const}\NormalTok{ mean }\OperatorTok{=}\NormalTok{ values}\OperatorTok{.}\FunctionTok{reduce}\NormalTok{((sum}\OperatorTok{,}\NormalTok{ value) }\KeywordTok{=\textgreater{}}\NormalTok{ sum }\OperatorTok{+}\NormalTok{ value}\OperatorTok{,} \DecValTok{0}\NormalTok{) }\OperatorTok{/}\NormalTok{ values}\OperatorTok{.}\AttributeTok{length}\OperatorTok{;}
  \KeywordTok{const}\NormalTok{ variance }\OperatorTok{=}\NormalTok{ values}\OperatorTok{.}\FunctionTok{reduce}\NormalTok{(}
\NormalTok{    (sum}\OperatorTok{,}\NormalTok{ value) }\KeywordTok{=\textgreater{}}\NormalTok{ sum }\OperatorTok{+}\NormalTok{ ((value }\OperatorTok{{-}}\NormalTok{ mean) }\OperatorTok{**} \DecValTok{2}\NormalTok{)}\OperatorTok{,}
    \DecValTok{0}
\NormalTok{  ) }\OperatorTok{/}\NormalTok{ values}\OperatorTok{.}\AttributeTok{length}\OperatorTok{;}
  \ControlFlowTok{return}\NormalTok{ \{ mean}\OperatorTok{,} \DataTypeTok{stdDev}\OperatorTok{:} \BuiltInTok{Math}\OperatorTok{.}\FunctionTok{sqrt}\NormalTok{(variance) \}}\OperatorTok{;}
\NormalTok{\}}
\end{Highlighting}
\end{Shaded}

Code dùng population variance, chia cho N. Window 30 mẫu được coi là
toàn bộ baseline hiện tại, không phải mẫu dùng để suy rộng thống kê học
thuật.

\subsection{3.3 Evaluate mẫu mới}\label{33-evaluate-mux1eabu-mux1edbi}

\begin{Shaded}
\begin{Highlighting}[]
\ControlFlowTok{if}\NormalTok{ (}\OperatorTok{!}\BuiltInTok{Object}\OperatorTok{.}\FunctionTok{hasOwn}\NormalTok{(MIN\_DELTA}\OperatorTok{,}\NormalTok{ sensorType)) \{}
  \ControlFlowTok{return}\NormalTok{ \{ }\DataTypeTok{ready}\OperatorTok{:} \KeywordTok{false}\OperatorTok{,} \DataTypeTok{zScore}\OperatorTok{:} \KeywordTok{null}\OperatorTok{,} \DataTypeTok{isAnomaly}\OperatorTok{:} \KeywordTok{false}\NormalTok{ \}}\OperatorTok{;}
\NormalTok{\}}
\end{Highlighting}
\end{Shaded}

Dòng này là hàng rào bảo đảm gas và motion không bị đưa vào Z-score.

\begin{Shaded}
\begin{Highlighting}[]
\ControlFlowTok{if}\NormalTok{ (buffer}\OperatorTok{.}\AttributeTok{length} \OperatorTok{===}\NormalTok{ WINDOW\_SIZE) \{}
  \KeywordTok{const}\NormalTok{ \{ mean}\OperatorTok{,}\NormalTok{ stdDev \} }\OperatorTok{=} \FunctionTok{meanAndStdDev}\NormalTok{(buffer)}\OperatorTok{;}
  \ControlFlowTok{if}\NormalTok{ (stdDev }\OperatorTok{===} \DecValTok{0}\NormalTok{) \{}
\NormalTok{    result }\OperatorTok{=}\NormalTok{ \{ }\DataTypeTok{ready}\OperatorTok{:} \KeywordTok{true}\OperatorTok{,} \DataTypeTok{zScore}\OperatorTok{:} \DecValTok{0}\OperatorTok{,}\NormalTok{ mean}\OperatorTok{,}\NormalTok{ stdDev}\OperatorTok{,} \DataTypeTok{isAnomaly}\OperatorTok{:} \KeywordTok{false}\NormalTok{ \}}\OperatorTok{;}
\NormalTok{  \} }\ControlFlowTok{else}\NormalTok{ \{}
    \KeywordTok{const}\NormalTok{ zScore }\OperatorTok{=}\NormalTok{ (value }\OperatorTok{{-}}\NormalTok{ mean) }\OperatorTok{/}\NormalTok{ stdDev}\OperatorTok{;}
    \KeywordTok{const}\NormalTok{ isAnomaly }\OperatorTok{=} \BuiltInTok{Math}\OperatorTok{.}\FunctionTok{abs}\NormalTok{(zScore) }\OperatorTok{\textgreater{}}\NormalTok{ systemThresholds}\OperatorTok{.}\AttributeTok{zScoreLimit}
      \OperatorTok{\&\&} \BuiltInTok{Math}\OperatorTok{.}\FunctionTok{abs}\NormalTok{(value }\OperatorTok{{-}}\NormalTok{ mean) }\OperatorTok{\textgreater{}}\NormalTok{ MIN\_DELTA[sensorType]}\OperatorTok{;}
\NormalTok{    result }\OperatorTok{=}\NormalTok{ \{ }\DataTypeTok{ready}\OperatorTok{:} \KeywordTok{true}\OperatorTok{,}\NormalTok{ zScore}\OperatorTok{,}\NormalTok{ mean}\OperatorTok{,}\NormalTok{ stdDev}\OperatorTok{,}\NormalTok{ isAnomaly \}}\OperatorTok{;}
\NormalTok{  \}}
\NormalTok{\}}

\NormalTok{buffer}\OperatorTok{.}\FunctionTok{push}\NormalTok{(value)}\OperatorTok{;}
\ControlFlowTok{if}\NormalTok{ (buffer}\OperatorTok{.}\AttributeTok{length} \OperatorTok{\textgreater{}}\NormalTok{ WINDOW\_SIZE) buffer}\OperatorTok{.}\FunctionTok{shift}\NormalTok{()}\OperatorTok{;}
\end{Highlighting}
\end{Shaded}

Thứ tự rất quan trọng: evaluate bằng 30 mẫu cũ, rồi mới roll mẫu mới vào
window.

\subsection{3.4 Bốn rule builder}\label{34-bux1ed1n-rule-builder}

Sau refactor, mỗi sensor có một hàm dễ mở và giải thích:

\begin{Shaded}
\begin{Highlighting}[]
\NormalTok{buildGasCandidate()}
\NormalTok{buildTemperatureCandidate()}
\NormalTok{buildMotionCandidate()}
\NormalTok{buildHumidityCandidate()}
\end{Highlighting}
\end{Shaded}

Hàm chính chỉ còn orchestration:

\begin{Shaded}
\begin{Highlighting}[]
\KeywordTok{const}\NormalTok{ context }\OperatorTok{=}\NormalTok{ \{ deviceId}\OperatorTok{,}\NormalTok{ timestamp}\OperatorTok{,}\NormalTok{ metrics \}}\OperatorTok{;}
\KeywordTok{const}\NormalTok{ temperatureStats }\OperatorTok{=} \FunctionTok{evaluateZScore}\NormalTok{(deviceId}\OperatorTok{,} \StringTok{\textquotesingle{}temperature\textquotesingle{}}\OperatorTok{,}\NormalTok{ metrics}\OperatorTok{.}\AttributeTok{temperature}\NormalTok{)}\OperatorTok{;}
\KeywordTok{const}\NormalTok{ humidityStats }\OperatorTok{=} \FunctionTok{evaluateZScore}\NormalTok{(deviceId}\OperatorTok{,} \StringTok{\textquotesingle{}humidity\textquotesingle{}}\OperatorTok{,}\NormalTok{ metrics}\OperatorTok{.}\AttributeTok{humidity}\NormalTok{)}\OperatorTok{;}
\KeywordTok{const}\NormalTok{ candidates }\OperatorTok{=}\NormalTok{ [}
  \FunctionTok{buildGasCandidate}\NormalTok{(context)}\OperatorTok{,}
  \FunctionTok{buildTemperatureCandidate}\NormalTok{(context}\OperatorTok{,}\NormalTok{ temperatureStats)}\OperatorTok{,}
  \FunctionTok{buildMotionCandidate}\NormalTok{(context)}\OperatorTok{,}
  \FunctionTok{buildHumidityCandidate}\NormalTok{(context}\OperatorTok{,}\NormalTok{ humidityStats)}
\NormalTok{]}\OperatorTok{.}\FunctionTok{filter}\NormalTok{(}\BuiltInTok{Boolean}\NormalTok{)}\OperatorTok{;}

\KeywordTok{const}\NormalTok{ created }\OperatorTok{=} \ControlFlowTok{await} \BuiltInTok{Promise}\OperatorTok{.}\FunctionTok{all}\NormalTok{(candidates}\OperatorTok{.}\FunctionTok{map}\NormalTok{(createAnomalyIfNew))}\OperatorTok{;}
\ControlFlowTok{return}\NormalTok{ created}\OperatorTok{.}\FunctionTok{filter}\NormalTok{(}\BuiltInTok{Boolean}\NormalTok{)}\OperatorTok{;}
\end{Highlighting}
\end{Shaded}

\texttt{candidate} là object trung gian mô tả anomaly và hành động cần
làm. Chỉ sau dedupe và action handling nó mới thành MongoDB document.

\subsection{3.5 Temperature rule}\label{35-temperature-rule}

\begin{Shaded}
\begin{Highlighting}[]
\KeywordTok{const}\NormalTok{ aboveMaximum }\OperatorTok{=}\NormalTok{ metrics}\OperatorTok{.}\AttributeTok{temperature} \OperatorTok{\textgreater{}}\NormalTok{ systemThresholds}\OperatorTok{.}\AttributeTok{tempMax}\OperatorTok{;}
\KeywordTok{const}\NormalTok{ belowMinimum }\OperatorTok{=}\NormalTok{ metrics}\OperatorTok{.}\AttributeTok{temperature} \OperatorTok{\textless{}}\NormalTok{ systemThresholds}\OperatorTok{.}\AttributeTok{tempMin}\OperatorTok{;}
\end{Highlighting}
\end{Shaded}

Hard bounds được ưu tiên. Nếu trong bounds và automatic mode bật, code
mới xét Z-score.

\begin{Shaded}
\begin{Highlighting}[]
\NormalTok{autoFan}\OperatorTok{:}\NormalTok{ metrics}\OperatorTok{.}\AttributeTok{temperature} \OperatorTok{\textgreater{}}\NormalTok{ statistics}\OperatorTok{.}\AttributeTok{mean}
\end{Highlighting}
\end{Shaded}

Z-score có thể âm hoặc dương. Fan chỉ bật khi bất thường theo hướng nóng
lên, không bật khi nhiệt độ giảm mạnh.

\subsection{3.6 Dedupe và action}\label{36-dedupe-vuxe0-action}

\begin{Shaded}
\begin{Highlighting}[]
\KeywordTok{function} \FunctionTok{dedupeKey}\NormalTok{(candidate) \{}
  \ControlFlowTok{return} \VerbatimStringTok{\textasciigrave{}}\SpecialCharTok{$\{}\NormalTok{candidate}\OperatorTok{.}\AttributeTok{deviceId}\SpecialCharTok{\}}\VerbatimStringTok{:}\SpecialCharTok{$\{}\NormalTok{candidate}\OperatorTok{.}\AttributeTok{sensorType}\SpecialCharTok{\}}\VerbatimStringTok{:}\SpecialCharTok{$\{}\NormalTok{candidate}\OperatorTok{.}\AttributeTok{triggerType}\SpecialCharTok{\}}\VerbatimStringTok{\textasciigrave{}}\OperatorTok{;}
\NormalTok{\}}
\end{Highlighting}
\end{Shaded}

\begin{Shaded}
\begin{Highlighting}[]
\KeywordTok{const}\NormalTok{ lastCreatedAt }\OperatorTok{=}\NormalTok{ anomalyDedupMap}\OperatorTok{.}\FunctionTok{get}\NormalTok{(key)}\OperatorTok{;}
\ControlFlowTok{if}\NormalTok{ (lastCreatedAt }\OperatorTok{!==} \KeywordTok{undefined} \OperatorTok{\&\&}\NormalTok{ now }\OperatorTok{{-}}\NormalTok{ lastCreatedAt }\OperatorTok{\textless{}}\NormalTok{ DEDUPE\_WINDOW\_MS) \{}
  \ControlFlowTok{return} \KeywordTok{null}\OperatorTok{;}
\NormalTok{\}}
\end{Highlighting}
\end{Shaded}

Key được reserve trước các thao tác async để hai message đến gần nhau
không cùng vượt qua kiểm tra.

Action được bọc \nolinkurl{try/catch} riêng. Nếu bật fan hoặc gửi email
thất bại, anomaly vẫn được ghi và \texttt{actionsTaken} chứa lý do lỗi.
Đây là graceful degradation.

\section{4. NC-03 - Time Series
ingestion}\label{4-nc-03---time-series-ingestion}

\subsection{4.1 Schema}\label{41-schema}

File:
\href{../backend/models/Schemas.js}{\nolinkurl{backend/models/Schemas.js}}

\begin{Shaded}
\begin{Highlighting}[]
\KeywordTok{const}\NormalTok{ SensorReadingSchema }\OperatorTok{=} \KeywordTok{new} \FunctionTok{Schema}\NormalTok{(\{}
  \DataTypeTok{timestamp}\OperatorTok{:}\NormalTok{ \{ }\DataTypeTok{type}\OperatorTok{:} \BuiltInTok{Date}\OperatorTok{,} \DataTypeTok{required}\OperatorTok{:} \KeywordTok{true}\OperatorTok{,} \ControlFlowTok{default}\OperatorTok{:} \BuiltInTok{Date}\OperatorTok{.}\AttributeTok{now}\NormalTok{ \}}\OperatorTok{,}
  \DataTypeTok{deviceId}\OperatorTok{:}\NormalTok{ \{ }\DataTypeTok{type}\OperatorTok{:} \BuiltInTok{String}\OperatorTok{,} \DataTypeTok{required}\OperatorTok{:} \KeywordTok{true}\OperatorTok{,} \DataTypeTok{index}\OperatorTok{:} \KeywordTok{true}\NormalTok{ \}}\OperatorTok{,}
  \DataTypeTok{metrics}\OperatorTok{:}\NormalTok{ \{}
    \DataTypeTok{temperature}\OperatorTok{:}\NormalTok{ \{ }\DataTypeTok{type}\OperatorTok{:} \BuiltInTok{Number}\OperatorTok{,} \DataTypeTok{required}\OperatorTok{:} \KeywordTok{true}\NormalTok{ \}}\OperatorTok{,}
    \DataTypeTok{humidity}\OperatorTok{:}\NormalTok{ \{ }\DataTypeTok{type}\OperatorTok{:} \BuiltInTok{Number}\OperatorTok{,} \DataTypeTok{required}\OperatorTok{:} \KeywordTok{true}\NormalTok{ \}}\OperatorTok{,}
    \DataTypeTok{gasLevel}\OperatorTok{:}\NormalTok{ \{ }\DataTypeTok{type}\OperatorTok{:} \BuiltInTok{Number}\OperatorTok{,} \DataTypeTok{required}\OperatorTok{:} \KeywordTok{true}\NormalTok{ \}}\OperatorTok{,}
    \DataTypeTok{lightIntensity}\OperatorTok{:}\NormalTok{ \{ }\DataTypeTok{type}\OperatorTok{:} \BuiltInTok{Number}\OperatorTok{,} \ControlFlowTok{default}\OperatorTok{:} \KeywordTok{null}\NormalTok{ \}}\OperatorTok{,}
    \DataTypeTok{motionDetected}\OperatorTok{:}\NormalTok{ \{ }\DataTypeTok{type}\OperatorTok{:} \BuiltInTok{Boolean}\OperatorTok{,} \DataTypeTok{required}\OperatorTok{:} \KeywordTok{true}\NormalTok{ \}}
\NormalTok{  \}}\OperatorTok{,}
  \DataTypeTok{actuatorStates}\OperatorTok{:}\NormalTok{ \{}
    \DataTypeTok{fanOn}\OperatorTok{:}\NormalTok{ \{ }\DataTypeTok{type}\OperatorTok{:} \BuiltInTok{Boolean}\OperatorTok{,} \DataTypeTok{required}\OperatorTok{:} \KeywordTok{true}\OperatorTok{,} \ControlFlowTok{default}\OperatorTok{:} \KeywordTok{false}\NormalTok{ \}}\OperatorTok{,}
    \DataTypeTok{buzzerActive}\OperatorTok{:}\NormalTok{ \{ }\DataTypeTok{type}\OperatorTok{:} \BuiltInTok{Boolean}\OperatorTok{,} \DataTypeTok{required}\OperatorTok{:} \KeywordTok{true}\OperatorTok{,} \ControlFlowTok{default}\OperatorTok{:} \KeywordTok{false}\NormalTok{ \}}\OperatorTok{,}
    \DataTypeTok{alertLedActive}\OperatorTok{:}\NormalTok{ \{ }\DataTypeTok{type}\OperatorTok{:} \BuiltInTok{Boolean}\OperatorTok{,} \DataTypeTok{required}\OperatorTok{:} \KeywordTok{true}\OperatorTok{,} \ControlFlowTok{default}\OperatorTok{:} \KeywordTok{false}\NormalTok{ \}}
\NormalTok{  \}}
\NormalTok{\}}\OperatorTok{,}\NormalTok{ \{}
  \DataTypeTok{autoCreate}\OperatorTok{:} \KeywordTok{false}\OperatorTok{,}
  \DataTypeTok{autoIndex}\OperatorTok{:} \KeywordTok{false}
\NormalTok{\})}\OperatorTok{;}
\end{Highlighting}
\end{Shaded}

\texttt{autoCreate=false} ngăn Mongoose tạo collection thường trước khi
service có cơ hội tạo Time Series collection.

\subsection{4.2 Tạo và verify
collection}\label{42-tux1ea1o-vuxe0-verify-collection}

File:
\href{../backend/services/timeSeriesService.js}{\nolinkurl{backend/services/timeSeriesService.js}}

\begin{Shaded}
\begin{Highlighting}[]
\KeywordTok{const}\NormalTok{ EXPIRE\_AFTER\_SECONDS }\OperatorTok{=} \DecValTok{30} \OperatorTok{*} \DecValTok{24} \OperatorTok{*} \DecValTok{60} \OperatorTok{*} \DecValTok{60}\OperatorTok{;}
\KeywordTok{const}\NormalTok{ TIME\_SERIES\_OPTIONS }\OperatorTok{=}\NormalTok{ \{}
  \DataTypeTok{timeField}\OperatorTok{:} \StringTok{\textquotesingle{}timestamp\textquotesingle{}}\OperatorTok{,}
  \DataTypeTok{metaField}\OperatorTok{:} \StringTok{\textquotesingle{}deviceId\textquotesingle{}}\OperatorTok{,}
  \DataTypeTok{granularity}\OperatorTok{:} \StringTok{\textquotesingle{}seconds\textquotesingle{}}
\NormalTok{\}}\OperatorTok{;}
\end{Highlighting}
\end{Shaded}

Hàm \textbf{ensureSensorReadingsTimeSeries}:

\begin{enumerate}
\def\labelenumi{\arabic{enumi}.}
\tightlist
\item
  Gọi \texttt{listCollections()} tìm \texttt{sensor\_readings}.
\item
  Nếu chưa có, gọi \texttt{createCollection()} với \texttt{timeseries}
  và TTL.
\item
  Đọc metadata lại.
\item
  \texttt{validateCollectionOptions()} so từng option.
\item
  Sai cấu hình thì throw để startup fail rõ ràng.
\end{enumerate}

\subsection{4.3 Validate và ingest}\label{43-validate-vuxe0-ingest}

File:
\href{../backend/services/telemetryService.js}{\nolinkurl{backend/services/telemetryService.js}}

\begin{Shaded}
\begin{Highlighting}[]
\KeywordTok{async} \KeywordTok{function} \FunctionTok{ingestTelemetry}\NormalTok{(payload}\OperatorTok{,}\NormalTok{ expectedDeviceId) \{}
  \KeywordTok{const}\NormalTok{ normalized }\OperatorTok{=} \FunctionTok{validateTelemetryPayload}\NormalTok{(payload}\OperatorTok{,}\NormalTok{ expectedDeviceId)}\OperatorTok{;}
  \KeywordTok{const}\NormalTok{ reading }\OperatorTok{=} \ControlFlowTok{await}\NormalTok{ SensorReading}\OperatorTok{.}\FunctionTok{create}\NormalTok{(normalized)}\OperatorTok{;}
  \KeywordTok{const}\NormalTok{ anomalies }\OperatorTok{=} \ControlFlowTok{await} \FunctionTok{processTelemetryForAnomalies}\NormalTok{(reading)}\OperatorTok{;}
  \ControlFlowTok{return}\NormalTok{ \{ reading}\OperatorTok{,}\NormalTok{ anomalies \}}\OperatorTok{;}
\NormalTok{\}}
\end{Highlighting}
\end{Shaded}

\texttt{expectedDeviceId} lấy từ MQTT topic. Nếu payload khai một
deviceId khác topic, validation từ chối để tránh ghi dữ liệu sai thiết
bị.

Validation còn kiểm tra:

\begin{itemize}
\tightlist
\item
  Nhiệt độ -40 đến 85°C.
\item
  Độ ẩm 0 đến 100\%.
\item
  Gas và light là số hữu hạn không âm.
\item
  Motion và actuator states là boolean.
\item
  Timestamp phải parse được.
\end{itemize}

\subsection{4.4 Historical route}\label{44-historical-route}

File:
\href{../backend/routes/telemetry.js}{\nolinkurl{backend/routes/telemetry.js}}

\begin{Shaded}
\begin{Highlighting}[]
\KeywordTok{let}\NormalTok{ databaseQuery }\OperatorTok{=}\NormalTok{ SensorReading}\OperatorTok{.}\FunctionTok{find}\NormalTok{(query)}
  \OperatorTok{.}\FunctionTok{sort}\NormalTok{(\{ }\DataTypeTok{timestamp}\OperatorTok{:} \OperatorTok{{-}}\DecValTok{1}\NormalTok{ \})}
  \OperatorTok{.}\FunctionTok{limit}\NormalTok{(parsedLimit)}\OperatorTok{;}

\KeywordTok{const}\NormalTok{ readings }\OperatorTok{=} \ControlFlowTok{await}\NormalTok{ databaseQuery}\OperatorTok{;}
\NormalTok{readings}\OperatorTok{.}\FunctionTok{reverse}\NormalTok{()}\OperatorTok{;}
\end{Highlighting}
\end{Shaded}

Hai bước này trả N mẫu mới nhất nhưng vẫn theo thứ tự thời gian phù hợp
để vẽ chart.

\section{5. NC-05 - Nodemailer và
throttle}\label{5-nc-05---nodemailer-vuxe0-throttle}

File:
\href{../backend/services/emailService.js}{\nolinkurl{backend/services/emailService.js}}

\subsection{5.1 Hằng số}\label{51-hux1eb1ng-sux1ed1}

\begin{Shaded}
\begin{Highlighting}[]
\KeywordTok{const}\NormalTok{ THROTTLE\_WINDOW\_MS }\OperatorTok{=} \DecValTok{15} \OperatorTok{*} \DecValTok{60} \OperatorTok{*} \DecValTok{1000}\OperatorTok{;}
\KeywordTok{const}\NormalTok{ ALERT\_SEVERITIES }\OperatorTok{=} \KeywordTok{new} \BuiltInTok{Set}\NormalTok{([}\StringTok{\textquotesingle{}HIGH\textquotesingle{}}\OperatorTok{,} \StringTok{\textquotesingle{}CRITICAL\textquotesingle{}}\NormalTok{])}\OperatorTok{;}
\KeywordTok{const}\NormalTok{ emailThrottleMap }\OperatorTok{=} \KeywordTok{new} \BuiltInTok{Map}\NormalTok{()}\OperatorTok{;}
\end{Highlighting}
\end{Shaded}

\subsection{5.2 Validate config}\label{52-validate-config}

\texttt{getEmailConfig()} chỉ cho bật email khi đủ SMTP host, user và
password. Port phải là integer từ 1 đến 65535. \texttt{SMTP\_SECURE} mặc
định true nếu port 465. \texttt{ALERT\_RECIPIENT\_EMAIL} không còn là
recipient chính; nó chỉ fallback khi database chưa có user.

\subsection{5.3 Lấy recipient từ user
profile}\label{53-lux1ea5y-recipient-tux1eeb-user-profile}

File:
\href{../backend/services/userNotificationService.js}{\nolinkurl{backend/services/userNotificationService.js}}

\begin{Shaded}
\begin{Highlighting}[]
\KeywordTok{const}\NormalTok{ users }\OperatorTok{=} \ControlFlowTok{await}\NormalTok{ User}\OperatorTok{.}\FunctionTok{find}\NormalTok{(\{\})}
  \OperatorTok{.}\FunctionTok{select}\NormalTok{(}\StringTok{\textquotesingle{}email settings.emailEnabled\textquotesingle{}}\NormalTok{)}
  \OperatorTok{.}\FunctionTok{lean}\NormalTok{()}\OperatorTok{;}

\KeywordTok{const}\NormalTok{ recipients }\OperatorTok{=}\NormalTok{ users}
  \OperatorTok{.}\FunctionTok{filter}\NormalTok{((user) }\KeywordTok{=\textgreater{}}\NormalTok{ user}\OperatorTok{.}\AttributeTok{settings}\OperatorTok{?.}\AttributeTok{emailEnabled} \OperatorTok{!==} \KeywordTok{false}\NormalTok{)}
  \OperatorTok{.}\FunctionTok{map}\NormalTok{((user) }\KeywordTok{=\textgreater{}} \FunctionTok{normalizeEmail}\NormalTok{(user}\OperatorTok{.}\AttributeTok{email}\NormalTok{))}\OperatorTok{;}
\end{Highlighting}
\end{Shaded}

Email nhận cảnh báo chính là \texttt{User.email} đã nhập khi đăng ký.
Settings lưu \texttt{settings.emailEnabled}. Khi đã có user nhưng tất cả
đều tắt, resolver trả danh sách rỗng và không dùng fallback, nhờ đó lựa
chọn trên web có hiệu lực thật.

\subsection{5.4 Gửi email}\label{54-gux1eedi-email}

\begin{Shaded}
\begin{Highlighting}[]
\KeywordTok{const}\NormalTok{ severity }\OperatorTok{=} \BuiltInTok{String}\NormalTok{(alert}\OperatorTok{.}\AttributeTok{severity} \OperatorTok{||} \StringTok{\textquotesingle{}\textquotesingle{}}\NormalTok{)}\OperatorTok{.}\FunctionTok{toUpperCase}\NormalTok{()}\OperatorTok{;}
\ControlFlowTok{if}\NormalTok{ (}\OperatorTok{!}\NormalTok{ALERT\_SEVERITIES}\OperatorTok{.}\FunctionTok{has}\NormalTok{(severity)) \{}
  \ControlFlowTok{return}\NormalTok{ \{ }\DataTypeTok{status}\OperatorTok{:} \StringTok{\textquotesingle{}ignored\textquotesingle{}}\OperatorTok{,} \DataTypeTok{reason}\OperatorTok{:} \StringTok{\textquotesingle{}severity\textquotesingle{}}\NormalTok{ \}}\OperatorTok{;}
\NormalTok{\}}
\end{Highlighting}
\end{Shaded}

MEDIUM và LOW dừng trước khi chạm SMTP.

\begin{Shaded}
\begin{Highlighting}[]
\KeywordTok{const}\NormalTok{ throttleKey }\OperatorTok{=} \VerbatimStringTok{\textasciigrave{}}\SpecialCharTok{$\{}\NormalTok{alert}\OperatorTok{.}\AttributeTok{deviceId}\SpecialCharTok{\}}\VerbatimStringTok{:}\SpecialCharTok{$\{}\NormalTok{alert}\OperatorTok{.}\AttributeTok{sensorType}\SpecialCharTok{\}}\VerbatimStringTok{\textasciigrave{}}\OperatorTok{;}
\KeywordTok{const}\NormalTok{ now }\OperatorTok{=} \FunctionTok{nowProvider}\NormalTok{()}\OperatorTok{;}
\KeywordTok{const}\NormalTok{ lastSent }\OperatorTok{=}\NormalTok{ emailThrottleMap}\OperatorTok{.}\FunctionTok{get}\NormalTok{(throttleKey)}\OperatorTok{;}
\ControlFlowTok{if}\NormalTok{ (lastSent }\OperatorTok{!==} \KeywordTok{undefined} \OperatorTok{\&\&}\NormalTok{ now }\OperatorTok{{-}}\NormalTok{ lastSent }\OperatorTok{\textless{}}\NormalTok{ THROTTLE\_WINDOW\_MS) \{}
  \ControlFlowTok{return}\NormalTok{ \{ }\DataTypeTok{status}\OperatorTok{:} \StringTok{\textquotesingle{}throttled\textquotesingle{}}\OperatorTok{,} \DataTypeTok{lastSentAt}\OperatorTok{:} \KeywordTok{new} \BuiltInTok{Date}\NormalTok{(lastSent) \}}\OperatorTok{;}
\NormalTok{\}}
\end{Highlighting}
\end{Shaded}

Sau đó mới build message và gửi:

\begin{Shaded}
\begin{Highlighting}[]
\KeywordTok{const}\NormalTok{ providedRecipients }\OperatorTok{=} \ControlFlowTok{await} \FunctionTok{recipientProvider}\NormalTok{(config)}\OperatorTok{;}
\ControlFlowTok{if}\NormalTok{ (providedRecipients}\OperatorTok{.}\AttributeTok{length} \OperatorTok{===} \DecValTok{0}\NormalTok{) \{}
  \ControlFlowTok{return}\NormalTok{ \{ }\DataTypeTok{status}\OperatorTok{:} \StringTok{\textquotesingle{}skipped\textquotesingle{}}\OperatorTok{,} \DataTypeTok{reason}\OperatorTok{:} \StringTok{\textquotesingle{}no\_recipients\textquotesingle{}}\NormalTok{ \}}\OperatorTok{;}
\NormalTok{\}}

\KeywordTok{const}\NormalTok{ result }\OperatorTok{=} \ControlFlowTok{await}\NormalTok{ activeTransporter}\OperatorTok{.}\FunctionTok{sendMail}\NormalTok{(\{}
  \DataTypeTok{from}\OperatorTok{:}\NormalTok{ config}\OperatorTok{.}\AttributeTok{from}\OperatorTok{,}
  \OperatorTok{...}\NormalTok{recipientHeaders}\OperatorTok{,}
  \OperatorTok{...}\NormalTok{message}
\NormalTok{\})}\OperatorTok{;}

\NormalTok{emailThrottleMap}\OperatorTok{.}\FunctionTok{set}\NormalTok{(throttleKey}\OperatorTok{,}\NormalTok{ now)}\OperatorTok{;}
\end{Highlighting}
\end{Shaded}

Map chỉ được cập nhật sau khi \texttt{sendMail} resolve thành công.

Một recipient dùng header \texttt{to}; nhiều recipient dùng \texttt{bcc}
để họ không nhìn thấy địa chỉ của nhau. Map throttle chỉ được cập nhật
sau khi \texttt{sendMail} resolve thành công.

\subsection{5.5 HTML safety}\label{55-html-safety}

Mọi giá trị động được đưa qua \texttt{escapeHtml()} trước khi chèn vào
HTML. Điều này ngăn dữ liệu cảm biến hoặc device ID chứa HTML tùy ý làm
hỏng email.

Timestamp vẫn được lưu và xử lý nội bộ theo UTC. Riêng template email
gọi \texttt{formatVietnamTimestamp()} với timezone
\nolinkurl{Asia/Ho_Chi_Minh}, nên người nhận thấy định dạng
\texttt{YYYY-MM-DD\ HH:mm:ss\ UTC+7} mà không làm thay đổi dữ liệu gốc.

\section{6. Test map}\label{6-test-map}

{\def\LTcaptype{none} % do not increment counter
\begin{longtable}[]{@{}
  >{\raggedright\arraybackslash}p{(\linewidth - 2\tabcolsep) * \real{0.3100}}
  >{\raggedright\arraybackslash}p{(\linewidth - 2\tabcolsep) * \real{0.6400}}@{}}
\toprule\noalign{}
\begin{minipage}[b]{\linewidth}\raggedright
File test
\end{minipage} & \begin{minipage}[b]{\linewidth}\raggedright
Chứng minh
\end{minipage} \\
\midrule\noalign{}
\endhead
\bottomrule\noalign{}
\endlastfoot
\texttt{actuatorStateService.test.js} & Mapping actuator và telemetry
reconciliation \\
\texttt{anomalyEngine.test.js} & N=30, per-device, zero variance, min
delta, hard rules, dedupe, integrated action \\
\texttt{emailService.test.js} & Severity filter, 15 phút, profile
recipients, opt-out, HTML escaping \\
\texttt{userNotificationService.test.js} & Chuẩn hóa, loại trùng và lọc
user đã tắt email \\
\texttt{userRoutes.test.js} & Settings lưu
\texttt{settings.emailEnabled} đúng field \\
\texttt{telemetryService.test.js} & Payload validation, topic/device
match, edge alert, actuator confirmation \\
\texttt{telemetryRoutes.test.js} & Latest-N query và chronological
response \\
\texttt{timeSeriesService.test.js} & Collection options và fail-fast \\
\nolinkurl{scripts/testCb03Mock.js} & REST fan ON/status/OFF có JWT
trong simulation mode \\
\end{longtable}
}

\section{6.1 MQ2 readiness contract cần biết khi giải thích
telemetry}\label{61-mq2-readiness-contract-cux1ea7n-biux1ebft-khi-giux1ea3i-thuxedch-telemetry}

Firmware không còn biến mẫu ADC lúc MQ2 chưa ổn định thành PPM. Driver
có năm trạng thái: \texttt{burn\_in}, \texttt{warming\_up},
\texttt{calibrating}, \texttt{ready}, \texttt{fault}. Chỉ \texttt{ready}
mới gửi \texttt{metrics.gasLevel} dạng số; các trạng thái khác gửi
\texttt{null} cùng \texttt{sensorStatus.gas} để backend vẫn lưu được
chẩn đoán.

\texttt{telemetryService.js} kiểm tra tính nhất quán hai chiều:
\texttt{ready=true} bắt buộc có PPM hữu hạn, còn chưa ready bắt buộc
\texttt{gasLevel=null}. \texttt{anomalyEngine.js} bỏ qua gas khi giá trị
không hữu hạn, nên dữ liệu burn-in không tạo anomaly/email giả. Firmware
cũ và mock không có \texttt{sensorStatus} vẫn tương thích và được xem là
legacy-ready.

Đường local safety độc lập với cloud: FreeRTOS task ưu tiên cao đọc MQ2
mỗi 100 ms. Sensor ready dùng ngưỡng 600 PPM; sensor đang chuẩn bị chỉ
dùng ngưỡng điện áp ADC gần bão hòa để bật còi/LED/OLED tại chỗ. Raw
fallback không publish MQTT vì nó không phải PPM đã hiệu chuẩn.

\section{7. Điều nên và không nên học
thuộc}\label{7-ux111iux1ec1u-nuxean-vuxe0-khuxf4ng-nuxean-hux1ecdc-thuux1ed9c}

Nên nhớ:

\begin{itemize}
\tightlist
\item
  Tên bốn luồng và thứ tự các tầng.
\item
  Công thức Z-score, N=30, threshold 3.0 và minimum delta.
\item
  Topic control và telemetry.
\item
  QoS 1 là at least once.
\item
  Time Series options và TTL 30 ngày.
\item
  Email key và window 15 phút.
\end{itemize}

Không cần học thuộc:

\begin{itemize}
\tightlist
\item
  Toàn bộ HTML email template.
\item
  Tất cả Tailwind class.
\item
  Từng dòng setup MQTT/TLS.
\item
  Giá trị ngẫu nhiên trong mock telemetry.
\end{itemize}

Khi không nhớ một chi tiết, mở đúng file và giải thích theo luồng dữ
liệu. Khả năng định vị code và lý giải quyết định quan trọng hơn việc
đọc thuộc source.

\chapter{Kịch bản demo phần Thái
Kiệt}\label{kux1ecbch-bux1ea3n-demo-phux1ea7n-thuxe1i-kiux1ec7t}

Mục tiêu của kịch bản là tách rõ ba mức bằng chứng:

\begin{enumerate}
\def\labelenumi{\arabic{enumi}.}
\tightlist
\item
  \textbf{Unit/integration simulation:} chạy được trên máy không cần
  ESP32, EMQX, Atlas hay SMTP thật.
\item
  \textbf{Build verification:} chứng minh frontend và firmware biên dịch
  được.
\item
  \textbf{Physical end-to-end:} chứng minh broker, database, email và
  relay thật hoạt động với latency yêu cầu.
\end{enumerate}

Không dùng kết quả mức 1 để tuyên bố mức 3 đã pass.

\section{1. Chuẩn bị trước
demo}\label{1-chuux1ea9n-bux1ecb-trux1b0ux1edbc-demo}

\subsection{Phần mềm}\label{phux1ea7n-mux1ec1m}

\begin{itemize}
\tightlist
\item
  Node.js 18 trở lên.
\item
  Dependencies đã cài bằng \texttt{npm\ ci} trong \texttt{backend} và
  \texttt{frontend}.
\item
  PlatformIO để build/flash firmware.
\end{itemize}

\subsection{Dịch vụ thật}\label{dux1ecbch-vux1ee5-thux1eadt}

\begin{itemize}
\tightlist
\item
  MongoDB Atlas connection string.
\item
  EMQX host, username, password và CA certificate.
\item
  SMTP account/app password và email nhận cảnh báo.
\end{itemize}

\subsection{Phần cứng CB-03}\label{phux1ea7n-cux1ee9ng-cb-03}

\begin{itemize}
\tightlist
\item
  ESP32 DevKit V1.
\item
  Relay HL-52S, signal nối GPIO 16.
\item
  Fan DC dùng nguồn phù hợp, chung ground với relay/ESP32 theo sơ đồ
  nhóm.
\item
  Không cấp nguồn fan trực tiếp từ GPIO ESP32.
\end{itemize}

\section{2. Demo nhanh không cần phần
cứng}\label{2-demo-nhanh-khuxf4ng-cux1ea7n-phux1ea7n-cux1ee9ng}

Mở PowerShell:

\begin{Shaded}
\begin{Highlighting}[]
\FunctionTok{cd}\NormalTok{ C}\OperatorTok{:}\NormalTok{\textbackslash{}Users\textbackslash{}ASUS\textbackslash{}OneDrive\textbackslash{}Documents\textbackslash{}ChatGPT\textbackslash{}Saferoom{-}IoT\textbackslash{}backend}
\NormalTok{npm run test}\OperatorTok{:}\NormalTok{kiet}
\end{Highlighting}
\end{Shaded}

Lệnh này chạy test tập trung cho CB-03, NC-01, NC-03, NC-05 và kết thúc
bằng mô phỏng fan.

Kết quả bắt buộc:

\begin{Shaded}
\begin{Highlighting}[]
\NormalTok{\# fail 0}
\NormalTok{CB{-}03 PASS: unauthorized access blocked; authenticated fan ON,}
\NormalTok{status, and fan OFF succeeded.}
\end{Highlighting}
\end{Shaded}

Nếu có test fail, không tiếp tục tuyên bố chức năng ổn. Đọc tên subtest
và sửa nguyên nhân trước.

\section{3. Show từng phần bằng test có mục
tiêu}\label{3-show-tux1eebng-phux1ea7n-bux1eb1ng-test-cuxf3-mux1ee5c-tiuxeau}

\subsection{3.1 CB-03}\label{31-cb-03}

\begin{Shaded}
\begin{Highlighting}[]
\NormalTok{npm run test}\OperatorTok{:}\NormalTok{cb03}
\end{Highlighting}
\end{Shaded}

Luồng được chứng minh:

\begin{Shaded}
\begin{Highlighting}[]
\NormalTok{request không JWT {-}\textgreater{} 401}
\NormalTok{fan ON {-}\textgreater{} API success {-}\textgreater{} mock actuator fanOn=true}
\NormalTok{GET status {-}\textgreater{} fanOn=true}
\NormalTok{fan OFF {-}\textgreater{} API success}
\end{Highlighting}
\end{Shaded}

Điểm phải nói:

\begin{quote}
Đây là simulation test cho REST, auth, MQTT abstraction và state. Nó
chưa đo relay thật.
\end{quote}

\subsection{3.2 NC-01}\label{32-nc-01}

\begin{Shaded}
\begin{Highlighting}[]
\NormalTok{node }\OperatorTok{{-}{-}}\NormalTok{test }\OperatorTok{{-}{-}}\NormalTok{test{-}name{-}pattern}\OperatorTok{=}\StringTok{"Z{-}score|hard{-}rule|gas or motion|deduplicates"}\NormalTok{ test}\OperatorTok{/}\NormalTok{anomalyEngine}\OperatorTok{.}\FunctionTok{test}\OperatorTok{.}\FunctionTok{js}
\end{Highlighting}
\end{Shaded}

Các bằng chứng cần chỉ:

\begin{itemize}
\tightlist
\item
  30 mẫu ổn định rồi spike.
\item
  Window tách theo device và metric.
\item
  Zero variance không chia cho 0.
\item
  Minimum delta loại nhiễu nhỏ.
\item
  Gas và motion không dùng Z-score.
\item
  Gas/motion/humidity đi đúng hard-rule matrix.
\item
  Dedupe loại anomaly lặp nhanh.
\end{itemize}

\subsection{3.3 NC-03}\label{33-nc-03}

\begin{Shaded}
\begin{Highlighting}[]
\NormalTok{node }\OperatorTok{{-}{-}}\NormalTok{test test}\OperatorTok{/}\NormalTok{timeSeriesService}\OperatorTok{.}\FunctionTok{test}\OperatorTok{.}\FunctionTok{js}\NormalTok{ test}\OperatorTok{/}\NormalTok{telemetryService}\OperatorTok{.}\FunctionTok{test}\OperatorTok{.}\FunctionTok{js}\NormalTok{ test}\OperatorTok{/}\NormalTok{telemetryRoutes}\OperatorTok{.}\FunctionTok{test}\OperatorTok{.}\FunctionTok{js}
\end{Highlighting}
\end{Shaded}

Các bằng chứng cần chỉ:

\begin{itemize}
\tightlist
\item
  Tạo \texttt{sensor\_readings} với Time Series options.
\item
  TTL 30 ngày.
\item
  Fail fast khi collection sai loại.
\item
  Validate telemetry và topic/device mismatch.
\item
  Historical query lấy N mẫu mới nhất và trả chronological order.
\end{itemize}

\subsection{3.4 NC-05}\label{34-nc-05}

\begin{Shaded}
\begin{Highlighting}[]
\NormalTok{node }\OperatorTok{{-}{-}}\NormalTok{test test}\OperatorTok{/}\NormalTok{emailService}\OperatorTok{.}\FunctionTok{test}\OperatorTok{.}\FunctionTok{js}
\end{Highlighting}
\end{Shaded}

Các bằng chứng cần chỉ:

\begin{itemize}
\tightlist
\item
  HIGH gửi email.
\item
  Cùng key bị throttle 15 phút.
\item
  Device hoặc sensor khác có window riêng.
\item
  MEDIUM không gọi SMTP.
\item
  Dynamic HTML được escape.
\end{itemize}

\subsection{3.5 Integrated NC-01 -\textgreater{} CB-03 -\textgreater{}
NC-05}\label{35-integrated-nc-01---cb-03---nc-05}

\begin{Shaded}
\begin{Highlighting}[]
\NormalTok{node }\OperatorTok{{-}{-}}\NormalTok{test }\OperatorTok{{-}{-}}\NormalTok{test{-}name{-}pattern}\OperatorTok{=}\StringTok{"runs the Z{-}score spike"}\NormalTok{ test}\OperatorTok{/}\NormalTok{anomalyEngine}\OperatorTok{.}\FunctionTok{test}\OperatorTok{.}\FunctionTok{js}
\end{Highlighting}
\end{Shaded}

Test này tạo 30 mẫu nhiệt độ ổn định rồi đưa spike 35°C. Nó kiểm tra:

\begin{Shaded}
\begin{Highlighting}[]
\NormalTok{Z{-}score \textgreater{} 3}
\NormalTok{{-}\textgreater{} tạo Z\_SCORE\_EXCEEDED / HIGH}
\NormalTok{{-}\textgreater{} publish fan ON}
\NormalTok{{-}\textgreater{} tạo anomaly document}
\NormalTok{{-}\textgreater{} gửi email lần đầu}
\NormalTok{{-}\textgreater{} lần gửi cùng key tiếp theo bị throttled}
\end{Highlighting}
\end{Shaded}

\section{4. Build verification}\label{4-build-verification}

\subsection{Frontend}\label{frontend}

\begin{Shaded}
\begin{Highlighting}[]
\FunctionTok{cd}\NormalTok{ C}\OperatorTok{:}\NormalTok{\textbackslash{}Users\textbackslash{}ASUS\textbackslash{}OneDrive\textbackslash{}Documents\textbackslash{}ChatGPT\textbackslash{}Saferoom{-}IoT\textbackslash{}frontend}
\NormalTok{npm run build}
\end{Highlighting}
\end{Shaded}

Kết quả bắt buộc là \texttt{built} và exit code 0. Cảnh báo bundle lớn
không phải compile failure nhưng nên được ghi nhận.

\subsection{Firmware}\label{firmware}

\begin{Shaded}
\begin{Highlighting}[]
\FunctionTok{cd}\NormalTok{ C}\OperatorTok{:}\NormalTok{\textbackslash{}Users\textbackslash{}ASUS\textbackslash{}OneDrive\textbackslash{}Documents\textbackslash{}ChatGPT\textbackslash{}Saferoom{-}IoT\textbackslash{}firmware}
\NormalTok{pio run}
\end{Highlighting}
\end{Shaded}

Kết quả bắt buộc là \texttt{{[}SUCCESS{]}}. Ghi lại phần trăm RAM và
flash để trả lời câu hỏi tài nguyên.

\section{5. Demo thật CB-03 /
TC-06}\label{5-demo-thux1eadt-cb-03--tc-06}

\subsection{5.1 Cấu hình}\label{51-cux1ea5u-huxecnh}

\begin{enumerate}
\def\labelenumi{\arabic{enumi}.}
\tightlist
\item
  Copy \nolinkurl{firmware/include/secrets.example.h} thành
  \nolinkurl{firmware/include/secrets.h}.
\item
  Điền WiFi, EMQX host, port 8883, username, password và CA certificate.
\item
  Trong \nolinkurl{backend/.env}, đặt:
\end{enumerate}

\begin{Shaded}
\begin{Highlighting}[]
\NormalTok{MOCK\_HARDWARE=false}
\NormalTok{DEVICE\_ID=safebox{-}01}
\NormalTok{MQTT\_BROKER\_URL=mqtts://\textless{}host\textgreater{}:8883}
\NormalTok{MQTT\_USERNAME=\textless{}username\textgreater{}}
\NormalTok{MQTT\_PASSWORD=\textless{}password\textgreater{}}
\end{Highlighting}
\end{Shaded}

\begin{enumerate}
\def\labelenumi{\arabic{enumi}.}
\setcounter{enumi}{3}
\tightlist
\item
  Bảo đảm backend và firmware dùng cùng \texttt{DEVICE\_ID}.
\end{enumerate}

\subsection{5.2 Chạy}\label{52-chux1ea1y}

Terminal 1:

\begin{Shaded}
\begin{Highlighting}[]
\FunctionTok{cd}\NormalTok{ firmware}
\NormalTok{pio run }\OperatorTok{{-}{-}}\NormalTok{target upload}
\NormalTok{pio device monitor}
\end{Highlighting}
\end{Shaded}

Terminal 2:

\begin{Shaded}
\begin{Highlighting}[]
\FunctionTok{cd}\NormalTok{ backend}
\NormalTok{npm }\FunctionTok{start}
\end{Highlighting}
\end{Shaded}

Terminal 3:

\begin{Shaded}
\begin{Highlighting}[]
\FunctionTok{cd}\NormalTok{ frontend}
\NormalTok{npm run dev}
\end{Highlighting}
\end{Shaded}

Mở \nolinkurl{http://localhost:3000}, đăng nhập và vào Overview.

\subsection{5.3 Bằng chứng cần
quay/chụp}\label{53-bux1eb1ng-chux1ee9ng-cux1ea7n-quaychux1ee5p}

\begin{enumerate}
\def\labelenumi{\arabic{enumi}.}
\tightlist
\item
  Dashboard hiển thị MQTT live và device online.
\item
  Nhấn fan ON.
\item
  Network tab cho thấy POST control trả 200 và
  \texttt{delivery=published}.
\item
  Serial Monitor in \texttt{{[}Fan{]}\ Relay\ state:\ ON}.
\item
  Relay click và fan quay.
\item
  Telemetry sau đó báo \texttt{fanOn=true}.
\item
  Lặp lại với OFF.
\end{enumerate}

\subsection{5.4 Đo latency}\label{54-ux111o-latency}

Mốc bắt đầu: thời điểm nhấn nút hoặc request rời browser.

Mốc kết thúc: thời điểm relay đổi trạng thái, có thể dùng video 60 fps
hoặc timestamp đồng bộ giữa frontend và Serial Monitor.

Tiêu chí TC-06:

\begin{Shaded}
\begin{Highlighting}[]
\NormalTok{latency \textless{} 500 ms}
\end{Highlighting}
\end{Shaded}

Ghi ít nhất 5 lần ON/OFF, báo median và worst case. Không chỉ đo một lần
đẹp nhất.

\section{6. Demo thật NC-03 /
TC-02}\label{6-demo-thux1eadt-nc-03--tc-02}

\subsection{6.1 Kiểm tra ingestion}\label{61-kiux1ec3m-tra-ingestion}

\begin{enumerate}
\def\labelenumi{\arabic{enumi}.}
\tightlist
\item
  Chạy ESP32, EMQX và backend thật.
\item
  Mở MongoDB Atlas Data Explorer.
\item
  Theo dõi \texttt{sensor\_readings} trong 20--30 giây.
\item
  Xác nhận timestamp cách nhau khoảng 2 giây và deviceId đúng.
\item
  Xác nhận metrics và actuatorStates có đủ field bắt buộc.
\end{enumerate}

\subsection{6.2 Kiểm tra collection
options}\label{62-kiux1ec3m-tra-collection-options}

Trong Atlas hoặc mongosh, kiểm tra metadata collection. Kết quả phải
tương đương:

\begin{Shaded}
\begin{Highlighting}[]
\FunctionTok{\{}
  \DataTypeTok{"timeseries"}\FunctionTok{:} \FunctionTok{\{}
    \DataTypeTok{"timeField"}\FunctionTok{:} \StringTok{"timestamp"}\FunctionTok{,}
    \DataTypeTok{"metaField"}\FunctionTok{:} \StringTok{"deviceId"}\FunctionTok{,}
    \DataTypeTok{"granularity"}\FunctionTok{:} \StringTok{"seconds"}
  \FunctionTok{\},}
  \DataTypeTok{"expireAfterSeconds"}\FunctionTok{:} \DecValTok{2592000}
\FunctionTok{\}}
\end{Highlighting}
\end{Shaded}

\subsection{6.3 Kiểm tra historical
API}\label{63-kiux1ec3m-tra-historical-api}

Sau khi đăng nhập, lấy JWT và gọi:

\begin{Shaded}
\begin{Highlighting}[]
\VariableTok{$token} \OperatorTok{=} \StringTok{\textquotesingle{}\textless{}JWT\_FROM\_LOGIN\textgreater{}\textquotesingle{}}
\VariableTok{$headers} \OperatorTok{=}\NormalTok{ @}\OperatorTok{\{}\NormalTok{ Authorization }\OperatorTok{=} \StringTok{"Bearer }\VariableTok{$token}\StringTok{"} \OperatorTok{\}}
\FunctionTok{Invoke{-}RestMethod}\NormalTok{ \textasciigrave{}}
  \OperatorTok{{-}}\NormalTok{Uri }\StringTok{\textquotesingle{}http://localhost:5000/api/v1/telemetry/historical?deviceId=safebox{-}01\&limit=10\textquotesingle{}}\NormalTok{ \textasciigrave{}}
  \OperatorTok{{-}}\NormalTok{Headers }\VariableTok{$headers}
\end{Highlighting}
\end{Shaded}

Xác nhận response có tối đa 10 mẫu mới nhất, timestamp tăng dần.

\section{7. Demo thật NC-01 /
TC-04}\label{7-demo-thux1eadt-nc-01--tc-04}

\subsection{Phương án an toàn}\label{phux1b0ux1a1ng-uxe1n-an-touxe0n}

Không làm nóng DHT22 quá mức hoặc tạo điều kiện nguy hiểm. Dùng
mock/test mode để trình bày thuật toán, hoặc dùng một script test
telemetry kiểm soát trong môi trường demo riêng.

Trình tự cần chứng minh:

\begin{enumerate}
\def\labelenumi{\arabic{enumi}.}
\tightlist
\item
  Bật automatic Z-score mode.
\item
  Nạp 30 mẫu ổn định quanh 25°C.
\item
  Nạp một mẫu spike đủ lớn nhưng vẫn hợp lệ, ví dụ 35°C.
\item
  Kiểm tra \texttt{anomaly\_logs} có:
\end{enumerate}

\begin{Shaded}
\begin{Highlighting}[]
\NormalTok{sensorType=temperature}
\NormalTok{triggerType=Z\_SCORE\_EXCEEDED}
\NormalTok{calculatedZScore \textgreater{} 3}
\NormalTok{severity=HIGH}
\end{Highlighting}
\end{Shaded}

\begin{enumerate}
\def\labelenumi{\arabic{enumi}.}
\setcounter{enumi}{4}
\tightlist
\item
  Kiểm tra Overview hiển thị anomaly.
\item
  Kiểm tra action fan ON được publish nếu spike hướng lên.
\end{enumerate}

Không chỉnh threshold gas local của ESP32 để giả lập Z-score; đây là hai
cơ chế khác nhau.

\section{8. Demo thật NC-05 /
TC-05}\label{8-demo-thux1eadt-nc-05--tc-05}

\subsection{8.1 Cấu hình SMTP}\label{81-cux1ea5u-huxecnh-smtp}

\begin{Shaded}
\begin{Highlighting}[]
\NormalTok{EMAIL\_ENABLED=true}
\NormalTok{SMTP\_HOST=smtp.gmail.com}
\NormalTok{SMTP\_PORT=587}
\NormalTok{SMTP\_SECURE=false}
\NormalTok{SMTP\_USER=\textless{}smtp{-}user\textgreater{}}
\NormalTok{SMTP\_PASS=\textless{}app{-}password\textgreater{}}
\NormalTok{SMTP\_FROM=SafeRoom IoT Box \textless{}smtp{-}user\textgreater{}}
\NormalTok{\# Chỉ dùng khi database chưa có user:}
\NormalTok{ALERT\_RECIPIENT\_EMAIL=\textless{}optional{-}fallback\textgreater{}}
\NormalTok{SMTP\_VERIFY\_ON\_STARTUP=true}
\end{Highlighting}
\end{Shaded}

Không đưa \texttt{.env} hoặc app password lên màn hình quay demo.

\subsection{8.2 Trình tự}\label{82-truxecnh-tux1ef1}

\begin{enumerate}
\def\labelenumi{\arabic{enumi}.}
\tightlist
\item
  Đăng nhập bằng tài khoản có email profile thật, mở Settings, bật
  \texttt{Nodemailer\ Email\ Alerts} và lưu.
\item
  Khởi động backend và xác nhận log
  \texttt{{[}Email{]}\ Alerts\ enabled.}.
\item
  Tạo một HIGH hoặc CRITICAL anomaly trong môi trường test an toàn.
\item
  Ghi timestamp anomaly.
\item
  Xác nhận email đến địa chỉ profile và nội dung đúng device, sensor,
  observed value, threshold.
\item
  Đo thời gian từ anomaly tới email; tiêu chí dưới 5 giây theo TC-05.
\item
  Tạo lại cùng \texttt{(deviceId,\ sensorType)} trong 15 phút.
\item
  Xác nhận không có email thứ hai và \texttt{actionsTaken} ghi
  \texttt{EMAIL\_THROTTLED}.
\item
  Tạo anomaly khác sensor hoặc device và xác nhận vẫn gửi được.
\end{enumerate}

\section{9. Bảng ghi kết quả vật
lý}\label{9-bux1ea3ng-ghi-kux1ebft-quux1ea3-vux1eadt-luxfd}

{\def\LTcaptype{none} % do not increment counter
\begin{longtable}[]{@{}
  >{\raggedright\arraybackslash}p{(\linewidth - 10\tabcolsep) * \real{0.1700}}
  >{\raggedleft\arraybackslash}p{(\linewidth - 10\tabcolsep) * \real{0.1100}}
  >{\raggedright\arraybackslash}p{(\linewidth - 10\tabcolsep) * \real{0.1700}}
  >{\raggedright\arraybackslash}p{(\linewidth - 10\tabcolsep) * \real{0.1700}}
  >{\raggedleft\arraybackslash}p{(\linewidth - 10\tabcolsep) * \real{0.1600}}
  >{\raggedright\arraybackslash}p{(\linewidth - 10\tabcolsep) * \real{0.1700}}@{}}
\toprule\noalign{}
\begin{minipage}[b]{\linewidth}\raggedright
Test
\end{minipage} & \begin{minipage}[b]{\linewidth}\raggedleft
Lần
\end{minipage} & \begin{minipage}[b]{\linewidth}\raggedright
Bắt đầu
\end{minipage} & \begin{minipage}[b]{\linewidth}\raggedright
Kết thúc
\end{minipage} & \begin{minipage}[b]{\linewidth}\raggedleft
Latency
\end{minipage} & \begin{minipage}[b]{\linewidth}\raggedright
Kết quả
\end{minipage} \\
\midrule\noalign{}
\endhead
\bottomrule\noalign{}
\endlastfoot
TC-06 Fan ON & 1 & & & & \\
TC-06 Fan OFF & 1 & & & & \\
TC-06 Fan ON & 2 & & & & \\
TC-06 Fan OFF & 2 & & & & \\
TC-05 Email & 1 & & & & \\
TC-05 Throttle & 1 & & & N/A & \\
\end{longtable}
}

\section{10. Cách xử lý khi demo
lỗi}\label{10-cuxe1ch-xux1eed-luxfd-khi-demo-lux1ed7i}

\subsection{API 401/403}\label{api-401403}

\begin{itemize}
\tightlist
\item
  Đăng nhập lại.
\item
  Kiểm tra localStorage có \texttt{safebox\_token}.
\item
  Không hardcode hoặc tắt auth để né lỗi.
\end{itemize}

\subsection{API 503 MQTT offline}\label{api-503-mqtt-offline}

\begin{itemize}
\tightlist
\item
  Kiểm tra \texttt{MOCK\_HARDWARE=false}.
\item
  Kiểm tra broker URL phải là \nolinkurl{mqtts://}.
\item
  Kiểm tra username/password và CA.
\item
  Kiểm tra log backend \nolinkurl{connected/subscribed}.
\end{itemize}

\subsection{Backend có MQTT nhưng device
offline}\label{backend-cuxf3-mqtt-nhux1b0ng-device-offline}

\begin{itemize}
\tightlist
\item
  Kiểm tra WiFi ESP32.
\item
  Kiểm tra device ID hai bên.
\item
  Kiểm tra status topic và retained online payload.
\end{itemize}

\subsection{Fan không quay nhưng Serial báo
ON}\label{fan-khuxf4ng-quay-nhux1b0ng-serial-buxe1o-on}

\begin{itemize}
\tightlist
\item
  Kiểm tra relay wiring, common ground và nguồn fan.
\item
  Đo GPIO 16.
\item
  Kiểm tra relay active HIGH hay active LOW thực tế. Code hiện giả định
  HIGH=ON theo pinout dự án.
\end{itemize}

\subsection{Không có email}\label{khuxf4ng-cuxf3-email}

\begin{itemize}
\tightlist
\item
  Kiểm tra \texttt{EMAIL\_ENABLED} và đủ SMTP variables.
\item
  Dùng app password, không dùng password tài khoản thường.
\item
  Kiểm tra severity phải HIGH/CRITICAL.
\item
  Kiểm tra có đang trong throttle window không.
\end{itemize}

\subsection{Không có dữ liệu
Atlas}\label{khuxf4ng-cuxf3-dux1eef-liux1ec7u-atlas}

\begin{itemize}
\tightlist
\item
  Kiểm tra MONGO\_URI và network access list.
\item
  Kiểm tra MQTT topic phải đúng \nolinkurl{safebox/{id}/telemetry}.
\item
  Kiểm tra backend log rejected payload do type/range/device mismatch.
\end{itemize}

\subsection{Gas cao hoặc không có
PPM}\label{gas-cao-houx1eb7c-khuxf4ng-cuxf3-ppm}

\begin{itemize}
\tightlist
\item
  Nếu web hiện \texttt{First\ burn-in}, giữ MQ2 cấp nguồn trong không
  khí sạch đủ 48 giờ; tiến độ được checkpoint vào NVS mỗi 10 phút.
\item
  Nếu hiện \texttt{Warming\ up}, chờ 5 phút; firmware đang dùng lại
  \texttt{Ro} đã lưu.
\item
  Nếu hiện \texttt{Calibrating}, không đưa cồn/khói/bật lửa gần cảm
  biến; firmware cần ít nhất 100 mẫu ổn định trong cửa sổ 60 giây.
\item
  Nếu hiện \texttt{Sensor\ fault}, đo AO sau cầu chia áp và kiểm tra
  đúng 15 kΩ / 10 kΩ, common ground, VCC 5 V.
\item
  TC-03 phải đo phản hồi local bằng phần cứng. Build pass không chứng
  minh thời gian thực tế dưới 200 ms.
\end{itemize}

\section{11. Bằng chứng dự
phòng}\label{11-bux1eb1ng-chux1ee9ng-dux1ef1-phuxf2ng}

Trước ngày vấn đáp nên chuẩn bị:

\begin{itemize}
\tightlist
\item
  Video TC-06 có đồng thời màn hình web, Serial Monitor và relay/fan.
\item
  Ảnh Atlas collection options và vài telemetry documents.
\item
  Email thật có timestamp và nội dung alert.
\item
  Terminal output \texttt{npm\ run\ test:kiet}.
\item
  Terminal output firmware/frontend build.
\end{itemize}

Bằng chứng quay trước là phương án dự phòng khi WiFi hoặc cloud lỗi,
không thay thế việc hiểu và giải thích code.

\end{document}
